% This file is a chapter fragment, in the same style as the other files under
% version/detailed/.  It is included by version/detailed.tex and must not
% contain a preamble of its own.
%
% Contents: the chapter which supplies the three remaining inputs of section
% 2.6 of the curve estimates -- the existence of a symmetric ample line bundle
% on an abelian variety, the positivity of the top self-intersection number of
% an ample line bundle, and the behaviour of intersection numbers under a
% finite surjective morphism -- and deduces from them, together with the
% formula [n]^*L = L^(n^2) of the previous chapter, that [n] is finite, flat
% and surjective for every n nonzero and that its degree is n^(2g).
%
% The intersection-number layer is built here from scratch, by the Snapper
% route: the Euler characteristic of a coherent sheaf is the one input taken on
% trust, and everything above it -- numerical polynomials, Snapper's theorem,
% multilinearity, symmetry, the projection formula, ampleness and positivity --
% is proved.
%
% Labels follow the project convention <type>:deg-<topic>.

\chapter{Ampleness, intersection numbers and the degree of $[n]$}
\label{chap:deg}

This chapter completes \Cref{sec:av-cube} of \Cref{chap:av-isogeny}.  What is
proved here is:

\begin{itemize}
  \item \emph{Intersection numbers.}  For line bundles $L_1,\dots,L_d$ on a
  projective scheme $X$ over an algebraically closed field, the integer
  $(L_1\cdots L_d)$ is defined (\Cref{def:deg-intersection}) as the coefficient
  of $m_1m_2\cdots m_d$ in the Snapper polynomial
  $(m_1,\dots,m_d)\mapsto\chi(X,L_1^{\otimes m_1}\otimes\cdots\otimes
  L_d^{\otimes m_d})$, and is proved to be symmetric
  (\Cref{lem:deg-intersection-symmetric}) and multilinear
  (\Cref{thm:deg-intersection-multilinear}).
  \item \emph{The projection formula},
  \Cref{prop:deg-projection-formula}: $(f^*L_1\cdots f^*L_d)=\deg
  f\cdot(L_1\cdots L_d)$ for a finite surjective morphism $f$ of projective
  integral schemes of dimension $d$.  This is
  \Cref{prop:av-degree-intersection}.
  \item \emph{Positivity}, \Cref{prop:deg-ample-positive}: $(L^g)>0$ for an
  ample line bundle $L$ on a projective scheme of dimension $g$.  This is
  \Cref{prop:av-ample-self-intersection}.
  \item \emph{A symmetric ample line bundle},
  \Cref{prop:deg-symmetric-ample}: every abelian variety carries one.  This is
  \Cref{prop:av-symmetric-ample}.
  \item \emph{The isogeny theorem}, \Cref{thm:deg-isogeny}: for every $n\neq0$
  the morphism $[n]$ is finite, flat and surjective.  This is
  \Cref{thm:av-isogeny-general}.
  \item \emph{The degree formula}, \Cref{thm:deg-degree-formula}:
  $\deg[n]=n^{2g}$.  This is \Cref{thm:av-degree-formula}.
\end{itemize}

The last two are the point of the chapter; the first four are what they are
proved from.  The one further input is
\Cref{thm:cube-mulby-pullback} of \Cref{chap:cube}, which for a symmetric line
bundle $L$ on an abelian variety gives $[n]^*L\cong L^{\otimes n^2}$; it is
used as a proved result and is not restated here.

Everything is built on \emph{one} unproved analytic input, isolated in
\Cref{sec:deg-euler}: the Euler characteristic $\chi(X,F)$ of a coherent sheaf
on a projective scheme over a field, that is the finite-dimensionality of the
cohomology groups, their vanishing above the dimension, and the additivity of
$\chi$ on short exact sequences.  Three further statements are left unproved
and each is marked where it occurs: the invariance of $\chi$ along a finite
morphism together with the projection formula for the direct image
(\Cref{prop:deg-euler-finite-morphism}), the projectivity of an abelian variety
(\Cref{prop:deg-abelian-projective}), and miracle flatness
(\Cref{prop:deg-miracle-flatness}).  Two more are quoted from the standard
theory of ample sheaves and of coherent sheaves on a noetherian scheme
(\Cref{lem:deg-ample-basic} and \Cref{lem:deg-devissage}).  Nothing else in
this chapter is asserted without proof.  \Cref{sec:deg-cost} counts what all of
it would cost to formalise, and is written so as to add to the count of
\Cref{sec:cube-cost} rather than to repeat it.

The route is Snapper's, as presented by Kleiman
\cite[\S1]{KleimanAmpleness}; the abelian-variety statements follow
\cite[\S6]{MumfordAbelianVarieties} and \cite[I, \S7]{MilneAbelianVarieties}.
Everything is written out in full, on the assumption that the reader has to
formalise from it and should not have to consult anything else.

\statusnote{The scope question recorded in \Cref{sec:av-cube} --- whether this
project builds the line-bundle and intersection-number layer of algebraic
geometry, or quotes $\deg[n]=n^{2g}$ as an external input --- has been settled
in favour of building it.  Nothing in this chapter is conditional on it.  What
is \emph{not} settled is when the Lean layer that \Cref{sec:deg-cost} and
\Cref{sec:cube-cost} price will be built; those two sections are written so
that the decision can be taken on their numbers.}

\section{Conventions}\label{sec:deg-conventions}

Throughout this chapter $k$ is an algebraically closed field; in particular it
is infinite, which is used in \Cref{lem:deg-general-hyperplane} and nowhere
else.  A \emph{variety} over $k$ is a $k$-scheme which is separated, of finite
type and geometrically integral over $k$, as in
\Cref{sec:cube-conventions}; in particular it is integral, irreducible,
connected and non-empty, and its closed points are exactly its $k$-rational
points and are dense in it.

All schemes occurring below are of finite type over $k$, hence noetherian, and
all sheaves of modules occurring below are coherent.  We write
$\mathbf P^N:=\mathbf P^N_k$ for projective $N$-space over $k$ and $\O(1)$ for
its Serre twisting sheaf.  A $k$-scheme is \emph{projective} if it admits a
closed immersion into some $\mathbf P^N$; a projective $k$-scheme is proper
over $k$.

For a line bundle --- that is, an invertible sheaf of $\O_X$-modules --- we
write $L^{\otimes m}$ for its $m$-th tensor power, $m\in\Z$, with
$L^{\otimes0}=\O_X$ and $L^{\otimes-1}=L^{-1}$ the inverse.  Given a tuple
$\underline L=(L_1,\dots,L_r)$ of line bundles on $X$ and
$\underline m=(m_1,\dots,m_r)\in\Z^r$ we abbreviate
\[
  \underline L^{\underline m}:=L_1^{\otimes m_1}\otimes\cdots\otimes
  L_r^{\otimes m_r}.
\]
For a coherent sheaf $F$ we write $\operatorname{Supp}F$ for its support, a
closed subset of $X$, and $\dim F:=\dim\operatorname{Supp}F$, with
$\dim\emptyset:=-1$ so that $\dim F=-1$ exactly when $F=0$.  For a morphism
$u\colon X'\to X$ and a sheaf $F$ on $X$ we write $u^*F$ for the pullback and,
for $u$ a closed immersion, $F|_{X'}$ for $u^*F$.

\begin{warning}[Composition order]\label{warn:deg-composition-order}
Composition is written in the classical order: $g\circ f$ means ``first $f$,
then $g$''.  This is the convention of \Cref{chap:av-isogeny} and
\Cref{chap:cube} --- see \Cref{warn:cube-composition-order} --- and it is the
opposite of the one used in the Lean sources, which write the same morphism as
\texttt{f} $\gg$ \texttt{g}.  Every composite in this chapter therefore
appears reversed in Lean.
\end{warning}

\begin{warning}[What the parallel layer supplies]
\label{warn:deg-parallel-layer}
As in \Cref{warn:cube-parallel-layer}, the tensor product of sheaves of modules
over a scheme, invertible sheaves and the Picard group $\mathrm{Pic}(X)$ are
being built in this project in parallel with this chapter, so they are treated
here as available and are \emph{not} counted as missing in
\Cref{sec:deg-cost}.  The same holds for the cohomology of a sheaf of modules
on a scheme and the direct images along a proper morphism: those are counted in
\Cref{sec:cube-cost}, as notions (N4) and (N6) there, and
\Cref{sec:deg-cost} charges only what sits on top of them.
\end{warning}

\section{The Euler characteristic}\label{sec:deg-euler}

Everything in this chapter rests on one number attached to a coherent sheaf on
a projective scheme over a field.  This section states what is needed about it,
and states it once, so that no later proof appeals to an unnamed cohomological
fact.

\begin{proposition}[Cohomology of a coherent sheaf on a projective scheme]
\label{prop:deg-cohomology-finiteness}
Let $X$ be a projective scheme over a field $k$ and let $F$ be a coherent sheaf
of $\O_X$-modules.  Then:
\begin{enumerate}
  \item each $H^i(X,F)$ is a finite-dimensional $k$-vector space;
  \item $H^i(X,F)=0$ for $i>\dim X$, and for $i<0$;
  \item a short exact sequence $0\to F'\to F\to F''\to0$ of coherent sheaves
  induces a long exact sequence
  \[
    0\to H^0(X,F')\to H^0(X,F)\to H^0(X,F'')\to H^1(X,F')\to\cdots
  \]
  of $k$-vector spaces;
  \item if $X$ is affine then $H^i(X,F)=0$ for $i>0$.
\end{enumerate}
\openobligation{AlgebraicGeometry: the cohomology of a coherent sheaf of
modules on a scheme, together with its finite dimensionality over the base
field for a proper scheme and its vanishing above the dimension. The
cohomology functors themselves are notion N4 of the formalisation cost section
of the chapter on the theorem of the cube; what is new here is the finiteness
and the vanishing.}
\end{proposition}
\begin{proof}
Not proved here.  Part (1) is \cite[III, \S5]{Hartshorne}, part (2) is
Grothendieck's vanishing theorem \cite[III, \S2]{Hartshorne}, part (3)
is the long exact sequence of derived functors, and part (4) is Serre's
vanishing theorem on an affine scheme, \cite[III, \S3]{Hartshorne}.
Parts (1) and (2) hold more generally for a proper scheme over a field; only
the projective case is used below.
\end{proof}

\begin{definition}[Euler characteristic]\label{def:deg-euler}
\uses{prop:deg-cohomology-finiteness}
Let $X$ be a projective scheme over a field $k$ and let $F$ be a coherent sheaf
of $\O_X$-modules.  The \emph{Euler characteristic} of $F$ is
\[
  \chi(X,F):=\sum_{i\geq0}(-1)^i\dim_kH^i(X,F)\;\in\;\Z .
\]
The sum is finite and each term is finite by
\Cref{prop:deg-cohomology-finiteness}(1) and (2), so $\chi(X,F)$ is a
well-defined integer.  We write $\chi(F)$ when $X$ is clear.
\end{definition}

\begin{lemma}[The Euler characteristic is additive]\label{lem:deg-euler-additive}
\uses{def:deg-euler}
Let $X$ be a projective scheme over a field $k$ and let
$0\to F'\to F\to F''\to0$ be a short exact sequence of coherent sheaves on
$X$.  Then
\[
  \chi(X,F)=\chi(X,F')+\chi(X,F'').
\]
In particular $\chi(X,0)=0$, and if $F'\cong F''$ then $\chi(F')=\chi(F'')$.
\end{lemma}
\begin{proof}
\uses{prop:deg-cohomology-finiteness, def:deg-euler}
By \Cref{prop:deg-cohomology-finiteness}(3) the sequence
\[
  0\to H^0(F')\to H^0(F)\to H^0(F'')\to H^1(F')\to H^1(F)\to\cdots
\]
is exact, and by parts (1) and (2) of the same proposition it consists of
finite-dimensional $k$-vector spaces and has only finitely many non-zero
terms.  For a bounded exact sequence $0\to V_0\to V_1\to\cdots\to V_n\to0$ of
finite-dimensional vector spaces one has $\sum_j(-1)^j\dim V_j=0$: this is the
rank--nullity theorem applied inductively, or the statement that the Euler
characteristic of an acyclic bounded complex vanishes.  Applying it to the
long exact sequence and grouping the terms in threes gives the assertion.
Invariance under isomorphism is the case $F''=0$ of the sequence
$0\to F'\to F\to0\to0$.
\end{proof}

\begin{lemma}[Euler characteristic of a finite scheme]
\label{lem:deg-euler-affine}
\uses{def:deg-euler}
Let $X$ be a projective scheme over $k$ which is affine, and let $F$ be
coherent on $X$.  Then $\chi(X,F)=\dim_k\Gamma(X,F)$.  In particular
$\chi(X,\O_X)=\dim_k\Gamma(X,\O_X)$, which is $>0$ if $X\neq\emptyset$.
\end{lemma}
\begin{proof}
\uses{prop:deg-cohomology-finiteness, def:deg-euler}
By \Cref{prop:deg-cohomology-finiteness}(4) all higher cohomology vanishes, so
only the $i=0$ term survives and $\chi(X,F)=\dim_kH^0(X,F)=\dim_k\Gamma(X,F)$.
For $F=\O_X$ and $X\neq\emptyset$ the ring $\Gamma(X,\O_X)$ is non-zero, so its
dimension over $k$ is at least one.
\end{proof}

\begin{proposition}[The Euler characteristic along a finite morphism]
\label{prop:deg-euler-finite-morphism}
Let $f\colon X\to Y$ be a finite morphism of projective schemes over $k$.
Then:
\begin{enumerate}
  \item for every coherent sheaf $F$ on $X$ the direct image $f_*F$ is coherent
  and $H^i(X,F)\cong H^i(Y,f_*F)$ for every $i$; in particular
  $\chi(X,F)=\chi(Y,f_*F)$;
  \item for every coherent $F$ on $X$ and every line bundle $N$ on $Y$ there
  is an isomorphism $f_*(F\otimes f^*N)\cong(f_*F)\otimes N$, natural in $F$
  and in $N$.
\end{enumerate}
\openobligation{AlgebraicGeometry: the invariance of cohomology under the
direct image along an affine morphism, and the projection formula for the
direct image. The projection formula is already charged as notion N5 of the
formalisation cost section of the chapter on the theorem of the cube; what is
new here is the affine invariance.}
\end{proposition}
\begin{proof}
Not proved here.  A finite morphism is affine, so (1) is
\cite[III, \S4 and \S8]{Hartshorne}: the higher direct images
of a quasi-coherent sheaf along an affine morphism vanish, whence the Leray
spectral sequence degenerates to the stated isomorphism; coherence of $f_*F$
is the finiteness of $f$.  Part (2) is the projection formula, which for $N$
locally free of finite rank is a local computation, \cite[II, \S5]{Hartshorne}.
\end{proof}

\begin{remark}[Why $\chi$ and not $h^0$]\label{rem:deg-why-chi}
The whole point of working with $\chi$ rather than with
$h^0=\dim_kH^0$ is \Cref{lem:deg-euler-additive}: $h^0$ is only left exact,
whereas $\chi$ is exactly additive, and the entire theory below is an exercise
in additivity.  The price is
\Cref{prop:deg-cohomology-finiteness}, that is, cohomology in \emph{all}
degrees, which \Cref{sec:cube-cost} already charges as notion (N4); the
chapter on the theorem of the cube needed degrees zero and one only, and this
one needs all of them.  There is no cheaper route: $\chi$ is what makes
$m\mapsto\chi(L^{\otimes m})$ polynomial, and $h^0$ is not polynomial in $m$
for small $m$.
\end{remark}

\section{Numerical polynomials}\label{sec:deg-numerical}

This section is pure combinatorics: no scheme occurs in it.  It is separated
out because it is the only part of the chapter that can be formalised today,
and because the whole definition of an intersection number is an application of
\Cref{lem:deg-difference-explicit}.

\begin{definition}[Numerical polynomial and the difference operators]
\label{def:deg-numerical}
Let $r\geq0$.  A function $P\colon\Z^r\to\Z$ is a \emph{numerical polynomial of
total degree at most $d$} if there is a polynomial
$\widetilde P\in\mathbf Q[T_1,\dots,T_r]$ of total degree at most $d$ with
$P(\underline m)=\widetilde P(\underline m)$ for every $\underline m\in\Z^r$.
For $1\leq j\leq r$ the \emph{$j$-th difference operator} is
\[
  (\Delta_jP)(\underline m):=P(\underline m+e_j)-P(\underline m),
\]
where $e_j$ is the $j$-th standard basis vector of $\Z^r$; more generally, for
$v\in\Z^r$ we write $(\Delta_vP)(\underline m):=P(\underline m+v)-P(\underline
m)$, so that $\Delta_j=\Delta_{e_j}$.  The operators $\Delta_v$ commute with
one another, being differences of the commuting shift operators
$(T_vP)(\underline m):=P(\underline m+v)$.
\end{definition}

\begin{lemma}[The representing polynomial is unique]
\label{lem:deg-numerical-unique}
\uses{def:deg-numerical}
If $\widetilde P,\widetilde Q\in\mathbf Q[T_1,\dots,T_r]$ agree at every point
of $\Z^r$ then $\widetilde P=\widetilde Q$.  Consequently a numerical
polynomial $P$ determines $\widetilde P$, hence its total degree and each of
its coefficients, unambiguously.
\end{lemma}
\begin{proof}
It suffices to show that a polynomial $R\in\mathbf Q[T_1,\dots,T_r]$ vanishing
on $\Z^r$ is zero, and this is an induction on $r$.  For $r=0$ there is
nothing to prove.  For $r\geq1$ write $R=\sum_{i}R_i(T_1,\dots,T_{r-1})T_r^i$.
For each fixed $\underline a\in\Z^{r-1}$ the one-variable polynomial
$\sum_iR_i(\underline a)T_r^i$ has infinitely many roots, namely every integer,
hence is zero; so every $R_i$ vanishes on $\Z^{r-1}$ and is zero by induction.
\end{proof}

\begin{lemma}[Differences lower the degree]
\label{lem:deg-difference-lowers-degree}
\uses{def:deg-numerical}
Let $P\colon\Z^r\to\Z$ be a numerical polynomial of total degree at most $d$
and let $v\in\Z^r$.  Then $\Delta_vP$ is a numerical polynomial of total degree
at most $d-1$; here ``of total degree at most $-1$'' means ``identically
zero''.  Consequently:
\begin{enumerate}
  \item $\Delta_{v_1}\cdots\Delta_{v_{d+1}}P=0$ for any $d+1$ vectors
  $v_1,\dots,v_{d+1}\in\Z^r$;
  \item $\Delta_{v_1}\cdots\Delta_{v_d}P$ is a constant function for any $d$
  vectors $v_1,\dots,v_d\in\Z^r$.
\end{enumerate}
\end{lemma}
\begin{proof}
\uses{lem:deg-numerical-unique}
Write $\widetilde P=\sum_{e\leq d}\widetilde P_e$ with $\widetilde P_e$
homogeneous of degree $e$.  Then
$\widetilde P(T+v)-\widetilde P(T)$ is a polynomial of total degree at most
$d$ whose degree-$d$ homogeneous part is
$\widetilde P_d(T+v)-\widetilde P_d(T)$; expanding
$\widetilde P_d(T+v)$ by multilinearity of the monomials, every term other
than $\widetilde P_d(T)$ involves at least one coordinate of $v$ and therefore
has degree at most $d-1$ in $T$.  So the degree-$d$ part of the difference
vanishes and $\Delta_vP$ is represented by a polynomial of total degree at
most $d-1$; it takes integer values because $P$ does.  If $d=0$ then
$\widetilde P$ is a constant and $\Delta_vP=0$.  Assertions (1) and (2) follow
by applying this $d+1$ respectively $d$ times, using
\Cref{lem:deg-numerical-unique} to speak of \emph{the} total degree at each
step; a numerical polynomial of total degree at most $0$ is a constant.
\end{proof}

\begin{lemma}[Vanishing differences characterise numerical polynomials]
\label{lem:deg-difference-criterion}
\uses{def:deg-numerical}
Let $P\colon\Z^r\to\Z$ and $d\geq-1$ be such that
$\Delta_{j_1}\Delta_{j_2}\cdots\Delta_{j_{d+1}}P=0$ for every choice of
indices $j_1,\dots,j_{d+1}\in\{1,\dots,r\}$, repetitions allowed.  Then $P$ is
a numerical polynomial of total degree at most $d$.
\end{lemma}
\begin{proof}
Induction on $r$, and inside it induction on $d$.

\emph{The case $r=1$.}  We must show: if $\Delta^{d+1}P=0$ then $P$ is a
numerical polynomial of degree at most $d$.  For $d=-1$ this says $P=0$.  Let
$d\geq0$ and assume the statement for $d-1$.  Then $\Delta^{d}(\Delta P)=0$, so
by the inductive hypothesis $\Delta P$ is a numerical polynomial of degree at
most $d-1$.  Set $c_i:=(\Delta^iP)(0)$ for $0\leq i\leq d$; these are integers.
We claim
\[
  P(m)=\sum_{i=0}^{d}c_i\binom{m}{i}\qquad(m\in\Z),
\]
where $\binom{m}{i}:=m(m-1)\cdots(m-i+1)/i!$, a polynomial of degree $i$ in $m$
taking integer values on $\Z$.  Indeed, put
$Q(m):=\sum_{i=0}^{d}c_i\binom{m}{i}$.  From the Pascal identity
$\binom{m+1}{i}-\binom{m}{i}=\binom{m}{i-1}$ one gets
$(\Delta Q)(m)=\sum_{i=1}^{d}c_i\binom{m}{i-1}$, and iterating,
$(\Delta^jQ)(0)=c_j$ for every $j\leq d$; also $\Delta^{d+1}Q=0$.  So
$R:=P-Q$ satisfies $\Delta^{d+1}R=0$ and $(\Delta^jR)(0)=0$ for every
$j\leq d$.  We claim that any such $R$ is identically zero, which gives $P=Q$
and finishes the case $r=1$.  The claim is an induction on $d$: for $d=-1$ it
says $R=0$; for $d\geq0$ the function $\Delta R$ satisfies the same two
hypotheses with $d$ replaced by $d-1$, since
$(\Delta^j(\Delta R))(0)=(\Delta^{j+1}R)(0)=0$ for $j\leq d-1$, hence is
identically zero by the inductive hypothesis; so $R$ is constant, and
$R(0)=c_0-c_0=0$.  Thus $P=Q$ is a numerical polynomial of degree at most
$d$.

\emph{The inductive step in $r$.}  Let $r\geq2$ and assume the statement for
$r-1$ variables and every $d$.  Fix $\underline m'=(m_2,\dots,m_r)$ and apply
the case $r=1$ to $m_1\mapsto P(m_1,\underline m')$, whose $(d+1)$-st
difference vanishes by hypothesis.  This gives
\[
  P(m_1,\underline m')=\sum_{i=0}^{d}c_i(\underline m')\binom{m_1}{i},
  \qquad c_i(\underline m'):=(\Delta_1^iP)(0,\underline m').
\]
Each $c_i\colon\Z^{r-1}\to\Z$ satisfies the hypothesis of the lemma with $d$
replaced by $d-i$: any $d-i+1$ difference operators in the variables
$m_2,\dots,m_r$ applied to $c_i$ give $\Delta_1^i$ followed by those
operators applied to $P$, evaluated at $m_1=0$, and that is a composite of
$d+1$ difference operators applied to $P$, hence zero.  By the inductive
hypothesis in $r$, each $c_i$ is a numerical polynomial of total degree at most
$d-i$.  Therefore $P$ is represented by
$\sum_{i\leq d}\widetilde c_i(T_2,\dots,T_r)\binom{T_1}{i}$, a polynomial of
total degree at most $\max_i\bigl((d-i)+i\bigr)=d$.
\end{proof}

\begin{lemma}[The top difference as an alternating sum]
\label{lem:deg-difference-explicit}
\uses{def:deg-numerical}
Let $P\colon\Z^r\to\Z$ be any function and let $v_1,\dots,v_d\in\Z^r$.  Then
\[
  (\Delta_{v_1}\cdots\Delta_{v_d}P)(\underline m)
  =\sum_{S\subseteq\{1,\dots,d\}}(-1)^{d-\#S}\,
   P\Bigl(\underline m+\sum_{i\in S}v_i\Bigr).
\]
Moreover, for $v,w\in\Z^r$ one has the operator identity
\[
  \Delta_{v+w}=\Delta_v+\Delta_w+\Delta_v\Delta_w .
\]
\end{lemma}
\begin{proof}
Write $T_v$ for the shift $(T_vP)(\underline m)=P(\underline m+v)$, so that
$\Delta_v=T_v-\mathrm{id}$ and $T_vT_w=T_{v+w}$.  All these operators commute.
Expanding the product $\prod_{i=1}^d(T_{v_i}-\mathrm{id})$ over the subsets $S$
of $\{1,\dots,d\}$ on which the factor $T_{v_i}$ is chosen gives
$\sum_S(-1)^{d-\#S}T_{\sum_{i\in S}v_i}$, which is the displayed formula.  For
the second identity,
\[
  \Delta_v+\Delta_w+\Delta_v\Delta_w
  =(T_v-\mathrm{id})+(T_w-\mathrm{id})+(T_v-\mathrm{id})(T_w-\mathrm{id})
  =T_vT_w-\mathrm{id}=T_{v+w}-\mathrm{id}=\Delta_{v+w}.
\]
\end{proof}

\section{Ample and very ample line bundles}\label{sec:deg-ample}

Ampleness is used twice: to produce, on an abelian variety, a line bundle whose
top self-intersection number is positive, and --- in the opposite direction ---
to force a scheme carrying an ample \emph{trivial} bundle to be finite.  This
section fixes the definition and collects the standard permanence properties;
they are quoted, not proved, and \Cref{sec:deg-cost} counts them.

\begin{definition}[Ample, very ample]\label{def:deg-ample}
Let $X$ be a scheme of finite type over $k$ and let $L$ be a line bundle on
$X$.
\begin{itemize}
  \item $L$ is \emph{ample} if for every coherent sheaf $F$ on $X$ there is an
  integer $n_0$ such that $F\otimes L^{\otimes n}$ is generated by its global
  sections for every $n\geq n_0$.
  \item $L$ is \emph{very ample} if there are $N\geq0$ and an immersion
  $i\colon X\to\mathbf P^N$ with $i^*\O(1)\cong L$.  If $X$ is proper over $k$
  such an immersion is automatically a closed immersion.
\end{itemize}
A \emph{symmetric} line bundle on a commutative group scheme $G$ over $k$ is a
line bundle $L$ with $[-1]^*L\cong L$, where $[-1]$ is the inversion morphism
of \Cref{def:cube-mulby-int}.
\end{definition}

\begin{lemma}[Permanence properties of ampleness]\label{lem:deg-ample-basic}
\uses{def:deg-ample}
Let $X$ be a proper scheme over $k$ and let $L,M$ be line bundles on $X$.
\begin{enumerate}
  \item If $L$ is ample and $M$ is generated by its global sections then
  $L\otimes M$ is ample.  In particular a tensor product of two ample line
  bundles is ample.
  \item For $n\geq1$: $L$ is ample if and only if $L^{\otimes n}$ is ample.
  \item If $u\colon X'\to X$ is a finite morphism of proper $k$-schemes ---
  in particular, if $u$ is a closed immersion or an isomorphism --- and $L$ is
  ample, then $u^*L$ is ample.
  \item $L$ is ample if and only if $L^{\otimes n}$ is very ample for some
  $n\geq1$.  Consequently a proper $k$-scheme carrying an ample line bundle is
  projective; and conversely $\O(1)$ is ample on $\mathbf P^N$, so by (3) every
  projective $k$-scheme carries an ample line bundle.
  \item If $L$ is ample then for every line bundle $M$ the bundle
  $M\otimes L^{\otimes n}$ is very ample for all $n\gg0$.
\end{enumerate}
\openobligation{AlgebraicGeometry: ampleness and very ampleness of an
invertible sheaf on a scheme, the notion of a sheaf generated by its global
sections, and the permanence properties of ampleness listed here. Mathlib
v4.33.1 has none of these.}
\end{lemma}
\begin{proof}
Not proved here; all five are classical.  See \cite[II, \S7]{Hartshorne}, where
(1), (2), (4) and (5) are proved for a scheme of finite type over a noetherian
ring, and \cite[\S1]{KleimanAmpleness}.  Part (3) for a closed immersion is
immediate from (4): compose the closed immersion with the one given by a very
ample power of $L$.  For a general finite $u$ it is the standard descent of
ampleness along a finite morphism, \cite[III, \S5]{Hartshorne}; only the case
of a closed immersion and the case of an isomorphism are used below.
\end{proof}

\begin{lemma}[A proper affine scheme is finite]
\label{lem:deg-proper-affine-finite}
Let $X$ be a $k$-scheme which is proper over $k$ and affine.  Then
$\Gamma(X,\O_X)$ is a finite-dimensional $k$-algebra, $X\cong\Spec
\Gamma(X,\O_X)$ is finite over $k$, and $\dim X\leq0$.
\end{lemma}
\begin{proof}
Since $X$ is affine, the counit of the adjunction between global sections and
$\Spec$ is an isomorphism $X\cong\Spec\Gamma(X,\O_X)$.  Since $X\to\Spec k$ is
universally closed and of finite type, $\Gamma(X,\O_X)$ is a finite
$k$-algebra; this is
\texttt{AlgebraicGeometry.finite\_appTop\_of\_universallyClosed}, which is in
Mathlib \texttt{v4.33.1} and is already listed among the available inputs in
\Cref{sec:cube-cost}.  A finite-dimensional $k$-algebra is artinian, so its
spectrum is a finite discrete space and has dimension at most $0$.
\end{proof}

\section{Snapper's theorem}\label{sec:deg-snapper}

\begin{lemma}[A hyperplane section avoiding the associated points]
\label{lem:deg-general-hyperplane}
Let $X$ be a projective scheme over the algebraically closed field $k$, let $L$
be a very ample line bundle on $X$ --- that is, $L\cong i^*\O(1)$ for a closed
immersion $i\colon X\to\mathbf P^N$ --- and let $F$ be a coherent sheaf on
$X$.  Then there is a global section $s\in\Gamma(X,L)$ such that
multiplication by $s$,
\[
  s\colon F\longrightarrow F\otimes L,
\]
is injective.  Its cokernel is $(F\otimes L)|_D$, where $D\subseteq X$ is the
closed subscheme cut out by $s$, and $\operatorname{Supp}\bigl((F\otimes
L)|_D\bigr)=\operatorname{Supp}F\cap D$ contains no generic point of an
irreducible component of $\operatorname{Supp}F$.
\end{lemma}
\begin{proof}
A noetherian scheme has only finitely many associated points --- the points
$x$ whose local ring has $\m_x$ among the associated primes of $F_x$ --- and
every generic point of an irreducible component of $\operatorname{Supp}F$ is
one of them.  Let $x_1,\dots,x_t$ be the associated points of $F$ and let
$i(x_1),\dots,i(x_t)$ be their images in $\mathbf P^N$.  A hyperplane
$H=V\bigl(\sum_ja_jT_j\bigr)$, $\underline a\in k^{N+1}\setminus\{0\}$,
contains $i(x_l)$ if and only if $\sum_ja_j\xi_j=0$ in the residue field
$\kappa(i(x_l))$, where $\underline\xi$ is a system of homogeneous coordinates
of $i(x_l)$; not all $\xi_j$ vanish, so this is a $k$-linear condition on
$\underline a$ cutting out a \emph{proper} subspace $W_l\subsetneq k^{N+1}$.
Since $k$ is infinite, $k^{N+1}$ is not the union of the finitely many proper
subspaces $W_1,\dots,W_t$; choose $\underline a$ outside all of them and let
$s\in\Gamma(X,L)$ be the image of the corresponding section
$\sum_ja_jT_j\in\Gamma(\mathbf P^N,\O(1))$ under
$\Gamma(\mathbf P^N,\O(1))\to\Gamma(X,i^*\O(1))=\Gamma(X,L)$.

By construction the zero scheme $D$ of $s$ contains no associated point of
$F$.  Injectivity of $s\colon F\to F\otimes L$ may be checked on stalks; at a
point $x$ the map is, after trivialising $L$ near $x$, multiplication by an
element $\sigma\in\O_{X,x}$ on $F_x$, and the kernel of that multiplication is
non-zero for some $x$ only if $\sigma$ lies in an associated prime of $F_x$,
which by the choice of $s$ happens for no $x$.  The cokernel is
$(F\otimes L)\otimes\O_X/(s)=(F\otimes L)|_D$ by definition of $D$, and its
support is $\operatorname{Supp}F\cap D$, which by construction misses the
generic points of the irreducible components of $\operatorname{Supp}F$.
\end{proof}

\begin{theorem}[Snapper]\label{thm:deg-snapper}
\uses{def:deg-euler, def:deg-numerical}
Let $X$ be a projective scheme over the algebraically closed field $k$, let $F$
be a coherent sheaf on $X$ and let $L_1,\dots,L_r$ be line bundles on $X$.
Then
\[
  P_{F,\underline L}\colon\Z^r\longrightarrow\Z,\qquad
  P_{F,\underline L}(\underline m):=
  \chi\bigl(X,F\otimes\underline L^{\underline m}\bigr)
\]
is a numerical polynomial of total degree at most $\dim F$.
\end{theorem}
\begin{proof}
\uses{lem:deg-euler-additive, lem:deg-difference-criterion,
lem:deg-difference-lowers-degree, lem:deg-general-hyperplane,
lem:deg-ample-basic}
Write $d:=\dim F$.  We argue by induction on $d$, the statement being proved
simultaneously for all $X$, all $r$, all tuples $\underline L$ and all
coherent $F$ with $\dim F\leq d$.

\emph{Base case $d=-1$.}  Then $F=0$, so $F\otimes\underline L^{\underline m}=0$
and $P_{F,\underline L}=0$ by \Cref{lem:deg-euler-additive}.

\emph{Reduction to very ample $L_j$.}  Fix an embedding $X\subseteq\mathbf P^N$
and put $A:=\O(1)|_X$, a very ample line bundle on $X$.  For each $j$ the sheaf
$L_j\otimes A^{\otimes n}$ is very ample for $n\gg0$ by
\Cref{lem:deg-ample-basic}(5), so we may write
$L_j\cong M_j\otimes N_j^{-1}$ with $M_j:=L_j\otimes A^{\otimes n_j}$ and
$N_j:=A^{\otimes n_j}$ both very ample.  Consider the tuple
$\underline M\underline N:=(M_1,\dots,M_r,N_1,\dots,N_r)$ of $2r$ very ample
line bundles, and suppose the theorem is known for it, so that
$Q:=P_{F,\underline M\underline N}$ is a numerical polynomial of total degree
at most $d$ on $\Z^{2r}$.  Then
\[
  P_{F,\underline L}(\underline m)
  =\chi\Bigl(F\otimes\bigotimes_jM_j^{\otimes m_j}\otimes
  \bigotimes_jN_j^{\otimes(-m_j)}\Bigr)
  =Q(m_1,\dots,m_r,-m_1,\dots,-m_r),
\]
and substituting the linear forms $m_j$ and $-m_j$ into a polynomial of total
degree at most $d$ yields a polynomial of total degree at most $d$.  So it is
enough to prove the theorem when every $L_j$ is very ample.

\emph{The inductive step.}  Let $d\geq0$ and assume the theorem for all
coherent sheaves of dimension at most $d-1$; let $F$ have $\dim F\leq d$ and
let $L_1,\dots,L_r$ all be very ample.  Fix $j$.  By
\Cref{lem:deg-general-hyperplane} applied to $F$ and $L_j$ there is a section
$s$ of $L_j$ and a short exact sequence
\[
  0\longrightarrow F\stackrel{s}{\longrightarrow}F\otimes L_j
  \longrightarrow G_j\longrightarrow0,\qquad
  G_j:=(F\otimes L_j)|_{D_j},
\]
with $\dim G_j\leq d-1$: indeed $\operatorname{Supp}G_j$ is a closed subset of
$\operatorname{Supp}F$ containing no generic point of an irreducible component
of the latter, hence meets each such component in a proper closed subset, of
strictly smaller dimension.  Tensoring the sequence with the invertible sheaf
$\underline L^{\underline m}$ is exact and gives, by
\Cref{lem:deg-euler-additive},
\[
  (\Delta_jP_{F,\underline L})(\underline m)
  =\chi\bigl(F\otimes\underline L^{\underline m}\otimes L_j\bigr)
  -\chi\bigl(F\otimes\underline L^{\underline m}\bigr)
  =\chi\bigl(G_j\otimes\underline L^{\underline m}\bigr)
  =P_{G_j,\underline L}(\underline m).
\]
By the inductive hypothesis $P_{G_j,\underline L}$ is a numerical polynomial of
total degree at most $d-1$, so by
\Cref{lem:deg-difference-lowers-degree}(1) any $d$ further difference operators
annihilate it.  Hence any $d+1$ difference operators annihilate
$P_{F,\underline L}$, and \Cref{lem:deg-difference-criterion} shows that
$P_{F,\underline L}$ is a numerical polynomial of total degree at most $d$.
\end{proof}

\begin{remark}[Proper rather than projective]\label{rem:deg-snapper-proper}
\Cref{thm:deg-snapper} holds for a proper scheme $X$ over $k$, with the same
statement; see \cite[\S1]{KleimanAmpleness} and the chapter \emph{Chow homology
and Chern classes} of \cite{stacks-project}.  The reduction to the projective
case is by Chow's lemma, \Cref{thm:chow-lemma} of \Cref{chap:chow-lemma}, in
the form \Cref{cor:chow-lemma-proper}, together with a d\'evissage on
$\dim F$.  It is not carried out here, and it is not needed: every scheme on
which this chapter forms an intersection number carries an ample line bundle
and is therefore projective (\Cref{lem:deg-ample-basic}(4)).  The two places
where this matters are \Cref{prop:deg-ample-positive}, whose hypothesis
already contains an ample bundle, and \Cref{lem:deg-torsion-finite}, where the
scheme is $A[n]$ and the ample bundle is the restriction of one on $A$.
\end{remark}

\section{Intersection numbers}\label{sec:deg-intersection}

\begin{definition}[Intersection number]\label{def:deg-intersection}
\uses{def:deg-euler}
Let $X$ be a projective scheme over $k$, let $d\geq0$, let $L_1,\dots,L_d$ be
line bundles on $X$ and let $F$ be a coherent sheaf on $X$.  The
\emph{intersection number of $L_1,\dots,L_d$ with coefficients in $F$} is
\[
  (L_1\cdots L_d\cdot F):=
  \sum_{S\subseteq\{1,\dots,d\}}(-1)^{d-\#S}\,
  \chi\Bigl(X,F\otimes\bigotimes_{i\in S}L_i\Bigr)\;\in\;\Z .
\]
We write $(L_1\cdots L_d):=(L_1\cdots L_d\cdot\O_X)$, and
$(L^d\cdot F):=(L\cdots L\cdot F)$ and $(L^d)$ when all the $L_i$ are equal to
$L$.  For $d=0$ the empty intersection number is $(\,\cdot\,F)=\chi(X,F)$.
\end{definition}

\begin{lemma}[The intersection number is the top difference of the Snapper
polynomial]\label{lem:deg-intersection-is-top-difference}
\uses{def:deg-intersection, thm:deg-snapper}
In the situation of \Cref{def:deg-intersection},
\[
  (L_1\cdots L_d\cdot F)=
  \bigl(\Delta_1\Delta_2\cdots\Delta_dP_{F,\underline L}\bigr)(0).
\]
If moreover $\dim F\leq d$, then $\Delta_1\cdots\Delta_dP_{F,\underline L}$ is
a constant function, so the same formula holds with $0$ replaced by any
$\underline m\in\Z^d$, and $(L_1\cdots L_d\cdot F)$ is the coefficient of the
monomial $T_1T_2\cdots T_d$ in the polynomial representing
$P_{F,\underline L}$.
\end{lemma}
\begin{proof}
\uses{lem:deg-difference-explicit, lem:deg-difference-lowers-degree,
thm:deg-snapper, def:deg-intersection}
By \Cref{lem:deg-difference-explicit} with $v_i=e_i$,
\[
  \bigl(\Delta_1\cdots\Delta_dP_{F,\underline L}\bigr)(0)
  =\sum_{S\subseteq\{1,\dots,d\}}(-1)^{d-\#S}P_{F,\underline L}
   \Bigl(\sum_{i\in S}e_i\Bigr)
  =\sum_S(-1)^{d-\#S}\chi\Bigl(F\otimes\bigotimes_{i\in S}L_i\Bigr),
\]
which is the definition of $(L_1\cdots L_d\cdot F)$.  Suppose now
$\dim F\leq d$.  By \Cref{thm:deg-snapper}, $P_{F,\underline L}$ is a numerical
polynomial of total degree at most $\dim F\leq d$, so by
\Cref{lem:deg-difference-lowers-degree}(2) the function
$\Delta_1\cdots\Delta_dP_{F,\underline L}$ is constant.  For the last
assertion, write the representing polynomial as
$\widetilde P=\sum_{\underline i}c_{\underline i}T^{\underline i}$, of total
degree at most $d$.  The operator $\Delta_1\cdots\Delta_d$ lowers the degree in
each variable $T_j$ by one, so it kills every monomial in which some $T_j$ does
not occur; a monomial in which every one of the $d$ variables occurs has total
degree at least $d$, hence exactly $d$, hence is $T_1T_2\cdots T_d$; and
$\Delta_1\cdots\Delta_d(T_1\cdots T_d)=1$.  So the constant equals
$c_{(1,\dots,1)}$.
\end{proof}

\begin{lemma}[Symmetry]\label{lem:deg-intersection-symmetric}
\uses{def:deg-intersection}
The intersection number $(L_1\cdots L_d\cdot F)$ is unchanged by permuting
$L_1,\dots,L_d$, and depends on each $L_i$ only through its isomorphism class.
\end{lemma}
\begin{proof}
\uses{def:deg-intersection, lem:deg-euler-additive}
The defining sum runs over the subsets of $\{1,\dots,d\}$ and its terms depend
only on the unordered family $\{L_i\}_{i\in S}$, so a permutation of the
indices merely permutes the terms.  Isomorphic line bundles give isomorphic
tensor products, hence equal Euler characteristics by
\Cref{lem:deg-euler-additive}.
\end{proof}

\begin{lemma}[Additivity in the coefficient sheaf]
\label{lem:deg-intersection-additive}
\uses{def:deg-intersection}
Let $0\to F'\to F\to F''\to0$ be a short exact sequence of coherent sheaves on
the projective $k$-scheme $X$ and let $L_1,\dots,L_d$ be line bundles.  Then
\[
  (L_1\cdots L_d\cdot F)=(L_1\cdots L_d\cdot F')+(L_1\cdots L_d\cdot F'').
\]
\end{lemma}
\begin{proof}
\uses{lem:deg-euler-additive, def:deg-intersection}
Tensoring the sequence with the invertible sheaf $\bigotimes_{i\in S}L_i$ is
exact, so \Cref{lem:deg-euler-additive} gives the identity term by term in the
defining sum.
\end{proof}

\begin{lemma}[Vanishing in small dimension]
\label{lem:deg-intersection-small-support}
\uses{def:deg-intersection}
If $\dim F<d$ then $(L_1\cdots L_d\cdot F)=0$.  In particular
$(L_1\cdots L_d\cdot F)=0$ whenever $F$ is supported on a closed subset of
dimension less than $d$, and $(L_1\cdots L_d)=0$ on a projective $k$-scheme of
dimension less than $d$.
\end{lemma}
\begin{proof}
\uses{thm:deg-snapper, lem:deg-difference-lowers-degree,
lem:deg-intersection-is-top-difference}
By \Cref{thm:deg-snapper} the function $P_{F,\underline L}$ is a numerical
polynomial of total degree at most $\dim F\leq d-1$, so by
\Cref{lem:deg-difference-lowers-degree}(1) the composite of the $d$ operators
$\Delta_1,\dots,\Delta_d$ annihilates it.  Now apply
\Cref{lem:deg-intersection-is-top-difference}.
\end{proof}

\begin{theorem}[Multilinearity]\label{thm:deg-intersection-multilinear}
\uses{def:deg-intersection}
Let $X$ be a projective scheme over $k$, let $F$ be coherent with $\dim F\leq
d$, let $M,N,L_2,\dots,L_d$ be line bundles on $X$ and let $n\in\Z$.  Then
\begin{enumerate}
  \item $(M\otimes N\cdot L_2\cdots L_d\cdot F)
        =(M\cdot L_2\cdots L_d\cdot F)+(N\cdot L_2\cdots L_d\cdot F)$;
  \item $(\O_X\cdot L_2\cdots L_d\cdot F)=0$;
  \item $(M^{\otimes n}\cdot L_2\cdots L_d\cdot F)
        =n\,(M\cdot L_2\cdots L_d\cdot F)$.
\end{enumerate}
By \Cref{lem:deg-intersection-symmetric} the same holds in every slot; thus
$(L_1\cdots L_d\cdot F)$ is a symmetric $\Z$-multilinear function of the
classes of $L_1,\dots,L_d$ in $\mathrm{Pic}(X)$.
\end{theorem}
\begin{proof}
\uses{thm:deg-snapper, lem:deg-difference-explicit,
lem:deg-difference-lowers-degree, lem:deg-intersection-is-top-difference,
lem:deg-intersection-symmetric}
Consider the tuple $\underline K:=(M,N,L_2,\dots,L_d)$ of $d+1$ line bundles,
with the variables of $\Z^{d+1}$ indexed by $0,1,2,\dots,d$, and put
$Q:=P_{F,\underline K}$, so that
$Q(m_0,m_1,m_2,\dots,m_d)=\chi\bigl(F\otimes M^{\otimes m_0}\otimes
N^{\otimes m_1}\otimes L_2^{\otimes m_2}\otimes\cdots\otimes
L_d^{\otimes m_d}\bigr)$.  By \Cref{thm:deg-snapper}, $Q$ is a numerical
polynomial of total degree at most $\dim F\leq d$.

\emph{(1).}  Write $v:=e_0+e_1\in\Z^{d+1}$.  Substituting $m_0=m_1$ into $Q$
gives exactly the Snapper polynomial of the tuple
$(M\otimes N,L_2,\dots,L_d)$, and the difference of that substituted function
in its first variable is $\Delta_v Q$ evaluated along the same substitution.
Iterating and using \Cref{lem:deg-intersection-is-top-difference} on both sides,
\[
  (M\otimes N\cdot L_2\cdots L_d\cdot F)
  =\bigl(\Delta_v\Delta_2\cdots\Delta_dQ\bigr)(0).
\]
By the operator identity of \Cref{lem:deg-difference-explicit},
$\Delta_v=\Delta_0+\Delta_1+\Delta_0\Delta_1$, so
\[
  \Delta_v\Delta_2\cdots\Delta_dQ
  =\Delta_0\Delta_2\cdots\Delta_dQ+\Delta_1\Delta_2\cdots\Delta_dQ
  +\Delta_0\Delta_1\Delta_2\cdots\Delta_dQ .
\]
The last summand is the composite of $d+1$ difference operators applied to a
numerical polynomial of total degree at most $d$, hence vanishes by
\Cref{lem:deg-difference-lowers-degree}(1).  The first two summands are
constant and equal, by \Cref{lem:deg-intersection-is-top-difference} applied to
the tuples $(M,L_2,\dots,L_d)$ and $(N,L_2,\dots,L_d)$, to
$(M\cdot L_2\cdots L_d\cdot F)$ and $(N\cdot L_2\cdots L_d\cdot F)$
respectively.  This is (1).

\emph{(2).}  The Snapper polynomial of the tuple $(\O_X,L_2,\dots,L_d)$ does
not depend on its first variable, since $\O_X^{\otimes m_1}=\O_X$ for every
$m_1$.  Hence $\Delta_1$ annihilates it, and
\Cref{lem:deg-intersection-is-top-difference} gives (2).

\emph{(3).}  For $n\geq0$ this follows from (1) and (2) by induction on $n$,
writing $M^{\otimes n}=M^{\otimes(n-1)}\otimes M$ and $M^{\otimes0}=\O_X$.  For
$n=-1$, apply (1) to $M$ and $M^{-1}$ and use (2):
$0=(\O_X\cdot\cdots)=(M\cdot\cdots)+(M^{-1}\cdot\cdots)$, so
$(M^{-1}\cdot\cdots)=-(M\cdot\cdots)$.  The case of general $n<0$ follows by
applying the non-negative case to $M^{-1}$.
\end{proof}

\begin{corollary}[Scaling the top self-intersection]
\label{cor:deg-intersection-powers}
\uses{thm:deg-intersection-multilinear}
Let $X$ be a projective $k$-scheme of dimension at most $g$, let $L$ be a line
bundle on $X$ and let $n\in\Z$.  Then
\[
  \bigl((L^{\otimes n})^g\bigr)=n^g\,(L^g).
\]
\end{corollary}
\begin{proof}
\uses{thm:deg-intersection-multilinear, lem:deg-intersection-symmetric}
Apply \Cref{thm:deg-intersection-multilinear}(3) in the first slot, then in the
second, and so on; by \Cref{lem:deg-intersection-symmetric} each application is
legitimate and each contributes a factor $n$.  Formally: induction on
$0\leq j\leq g$ shows
$\bigl((L^{\otimes n})^j\cdot L^{g-j}\bigr)=n^j(L^g)$, the inductive step being
\Cref{thm:deg-intersection-multilinear}(3) in the $(j+1)$-st slot.
\end{proof}

\begin{lemma}[D\'evissage of a coherent sheaf]\label{lem:deg-devissage}
Let $X$ be a noetherian scheme and let $F$ be a coherent sheaf on $X$.  Then
there is a finite filtration by coherent subsheaves
\[
  0=F_0\subseteq F_1\subseteq\cdots\subseteq F_t=F
\]
such that for each $j$ there are an integral closed subscheme
$\iota_j\colon Z_j\to X$ and a non-zero coherent sheaf of ideals
$\mathcal I_j\subseteq\O_{Z_j}$ with
$F_j/F_{j-1}\cong(\iota_j)_*\mathcal I_j$.
\openobligation{AlgebraicGeometry: the devissage of a coherent sheaf on a
noetherian scheme by pushforwards of ideal sheaves on integral closed
subschemes.}
\end{lemma}
\begin{proof}
Not proved here.  This is the standard d\'evissage of coherent sheaves; see the
section \emph{D\'evissage of coherent sheaves} of the chapter \emph{Cohomology
of schemes} in \cite{stacks-project}, and \cite[III, \S5]{Hartshorne}, where
the same filtration is the engine of the finiteness theorem for cohomology.
The construction is noetherian induction on $\operatorname{Supp}F$: one takes
$Z_1$ to be an irreducible component of $\operatorname{Supp}F$ of maximal
dimension with its reduced structure, and $F_1$ the subsheaf of sections
supported on $Z_1$ whose restriction to $Z_1$ has no embedded contribution.
\end{proof}

\begin{proposition}[The top intersection number sees only the generic rank]
\label{prop:deg-intersection-rank}
\uses{def:deg-intersection}
Let $Y$ be a projective integral $k$-scheme of dimension $d$ with generic point
$\eta$, let $L_1,\dots,L_d$ be line bundles on $Y$ and let $F$ be a coherent
sheaf on $Y$.  Put $r:=\dim_{\kappa(\eta)}F_\eta$, the \emph{generic rank} of
$F$.  Then
\[
  (L_1\cdots L_d\cdot F)=r\,(L_1\cdots L_d).
\]
\end{proposition}
\begin{proof}
\uses{lem:deg-devissage, lem:deg-intersection-additive,
lem:deg-intersection-small-support, def:deg-intersection}
Write $\phi(F):=(L_1\cdots L_d\cdot F)$.  By
\Cref{lem:deg-intersection-additive}, $\phi$ is additive on short exact
sequences; by \Cref{lem:deg-intersection-small-support}, $\phi(G)=0$ whenever
$\dim G<d$.

Take a filtration of $F$ as in \Cref{lem:deg-devissage}, with graded pieces
$(\iota_j)_*\mathcal I_j$ for integral closed subschemes $Z_j\subseteq Y$ and
non-zero ideal sheaves $\mathcal I_j\subseteq\O_{Z_j}$.  Additivity gives
$\phi(F)=\sum_j\phi\bigl((\iota_j)_*\mathcal I_j\bigr)$.

Since $Y$ is integral of dimension $d$, the only integral closed subscheme of
$Y$ of dimension $d$ is $Y$ itself: a proper closed subset of an irreducible
scheme of finite type over a field has strictly smaller dimension.  So for
each $j$ either $\dim Z_j<d$, in which case
$\dim\bigl((\iota_j)_*\mathcal I_j\bigr)\leq\dim Z_j<d$ and the term vanishes,
or $Z_j=Y$ and the term is $\phi(\mathcal I_j)$ for a non-zero ideal sheaf
$\mathcal I_j\subseteq\O_Y$.  In the latter case the exact sequence
$0\to\mathcal I_j\to\O_Y\to\O_Y/\mathcal I_j\to0$ has a third term supported on
the proper closed subset $V(\mathcal I_j)\subsetneq Y$, of dimension $<d$, so
$\phi(\mathcal I_j)=\phi(\O_Y)=(L_1\cdots L_d)$.

It remains to count the indices $j$ with $Z_j=Y$.  Taking stalks at $\eta$ is
exact, so
$\dim_{\kappa(\eta)}F_\eta=\sum_j\dim_{\kappa(\eta)}\bigl((\iota_j)_*\mathcal
I_j\bigr)_\eta$.  If $Z_j\neq Y$ then $\eta\notin Z_j$ and the stalk is $0$.
If $Z_j=Y$ then the stalk is $(\mathcal I_j)_\eta$, a non-zero ideal of the
field $\O_{Y,\eta}=\kappa(\eta)$, hence $\kappa(\eta)$ itself, of dimension
one.  So the number of $j$ with $Z_j=Y$ is exactly $r$, and
$\phi(F)=r\,(L_1\cdots L_d)$.
\end{proof}

\begin{proposition}[Projection formula for a finite surjective morphism]
\label{prop:deg-projection-formula}
\uses{def:deg-intersection, def:av-degree}
Let $f\colon X\to Y$ be a finite surjective morphism of projective integral
$k$-schemes of the same dimension $d$, and let $L_1,\dots,L_d$ be line bundles
on $Y$.  Then
\[
  \bigl(f^*L_1\cdots f^*L_d\bigr)_X=\deg f\cdot\bigl(L_1\cdots L_d\bigr)_Y,
\]
where $\deg f:=\dim_{\kappa(\eta_Y)}(f_*\O_X)_{\eta_Y}$ is the generic rank of
$f_*\O_X$, equal to the degree of the field extension
$\kappa(\eta_X)/\kappa(\eta_Y)$ of function fields.  If moreover $f$ is flat,
$f_*\O_X$ is locally free and $\deg f$ agrees with the degree of
\Cref{def:av-degree}.
\end{proposition}
\begin{proof}
\uses{prop:deg-euler-finite-morphism, prop:deg-intersection-rank,
def:deg-intersection, lem:deg-euler-additive, def:av-degree}
For every $S\subseteq\{1,\dots,d\}$ put $L_S:=\bigotimes_{i\in S}L_i$, so that
$\bigotimes_{i\in S}f^*L_i\cong f^*L_S$.  By
\Cref{prop:deg-euler-finite-morphism}(1) and (2),
\[
  \chi\bigl(X,f^*L_S\bigr)=\chi\bigl(Y,f_*f^*L_S\bigr)
  =\chi\bigl(Y,(f_*\O_X)\otimes L_S\bigr),
\]
the second equality because $f_*(\O_X\otimes f^*L_S)\cong(f_*\O_X)\otimes L_S$.
Summing over $S$ with the signs of \Cref{def:deg-intersection} gives
\[
  \bigl(f^*L_1\cdots f^*L_d\bigr)_X
  =\bigl(L_1\cdots L_d\cdot f_*\O_X\bigr)_Y .
\]
Now apply \Cref{prop:deg-intersection-rank} on $Y$: the right-hand side is
$r\,(L_1\cdots L_d)_Y$ with $r$ the generic rank of $f_*\O_X$.

Finally we identify $r$.  A finite morphism is affine, and $f_*\O_X$ is the
quasi-coherent sheaf of $\O_Y$-algebras whose stalk at $\eta_Y$ is
$(f_*\O_X)_{\eta_Y}=\Gamma\bigl(f^{-1}(\Spec R),\O_X\bigr)\otimes_R
\kappa(\eta_Y)$ for an affine open $\Spec R\subseteq Y$ containing $\eta_Y$;
since $f$ is finite, surjective and $X$, $Y$ are integral, $f$ is dominant and
the fibre of $f$ over $\eta_Y$ is the spectrum of the field
$\kappa(\eta_X)$, viewed as a $\kappa(\eta_Y)$-algebra.  Hence
$r=[\kappa(\eta_X):\kappa(\eta_Y)]$.  If $f$ is in addition flat then $f_*\O_X$
is locally free of finite rank and that rank, which is what
\Cref{def:av-degree} calls $\deg f$, is computed at any point, in particular at
$\eta_Y$, where it is $r$.
\end{proof}

\section{Positivity}\label{sec:deg-positivity}

\begin{proposition}[The top self-intersection of an ample bundle is positive]
\label{prop:deg-ample-positive}
\uses{def:deg-intersection, def:deg-ample}
Let $X$ be a proper scheme over $k$ of dimension $g\geq0$ and let $L$ be an
ample line bundle on $X$.  Then $X$ is projective and
\[
  (L^g)>0 .
\]
\end{proposition}
\begin{proof}
\uses{lem:deg-ample-basic, cor:deg-intersection-powers,
lem:deg-general-hyperplane, lem:deg-proper-affine-finite,
lem:deg-euler-additive, lem:deg-euler-affine, thm:deg-snapper,
lem:deg-difference-explicit, lem:deg-difference-lowers-degree,
lem:deg-intersection-is-top-difference, prop:deg-euler-finite-morphism}
That $X$ is projective is \Cref{lem:deg-ample-basic}(4).  We prove
$(L^g)>0$ by induction on $g$.

\emph{Reduction to $L$ very ample.}  By \Cref{lem:deg-ample-basic}(4) there is
$e\geq1$ with $L^{\otimes e}$ very ample, and by
\Cref{cor:deg-intersection-powers}
$\bigl((L^{\otimes e})^g\bigr)=e^g(L^g)$ with $e^g>0$.  So $(L^g)>0$ if and
only if $\bigl((L^{\otimes e})^g\bigr)>0$, and $L^{\otimes e}$ is again ample
by \Cref{lem:deg-ample-basic}(2).  We may therefore assume $L$ very ample, and
fix a closed immersion $X\subseteq\mathbf P^N$ with $L=\O(1)|_X$.

\emph{Base case $g=0$.}  Here $(L^0)=\chi(X,\O_X)$ by
\Cref{def:deg-intersection}.  A scheme of finite type over $k$ of dimension
$0$ is finite and discrete, hence affine, so
$\chi(X,\O_X)=\dim_k\Gamma(X,\O_X)>0$ by \Cref{lem:deg-euler-affine}; note
$X\neq\emptyset$, since $\dim\emptyset=-1$.

\emph{Inductive step.}  Let $g\geq1$ and assume the proposition in dimension
$g-1$.  Apply \Cref{lem:deg-general-hyperplane} with $F=\O_X$: there is
$s\in\Gamma(X,L)$, the restriction of a hyperplane form on $\mathbf P^N$, whose
zero scheme $D\subseteq X$ contains no generic point of an irreducible
component of $X$ and for which
\[
  0\longrightarrow\O_X\stackrel{s}{\longrightarrow}L
  \longrightarrow L|_D\longrightarrow0
\]
is exact.  Call this sequence the \emph{hyperplane sequence}.

\emph{The divisor $D$ has dimension exactly $g-1$.}  Let $X_1$ be an
irreducible component of $X$ of dimension $g$.  If $D\cap X_1$ were empty then
$X_1$ would be a closed subscheme of $\mathbf P^N\setminus H\cong\mathbf A^N$,
where $H$ is the hyperplane cutting out $s$; a closed subscheme of an affine
scheme is affine, and $X_1$ is proper over $k$, so
\Cref{lem:deg-proper-affine-finite} would give $\dim X_1\leq0$, contradicting
$\dim X_1=g\geq1$.  So $D\cap X_1\neq\emptyset$; it is cut out in $X_1$ by one
equation, so each of its irreducible components has codimension at most one in
$X_1$ by Krull's principal ideal theorem, whence $\dim(D\cap X_1)\geq g-1$.  On
the other hand $D$ contains no generic point of a component of $X$, so
$D\cap X_i\subsetneq X_i$ for every component $X_i$ and therefore
$\dim(D\cap X_i)<\dim X_i\leq g$.  Taking the maximum over the components gives
$\dim D=g-1$.

\emph{Restriction of $L$ to $D$.}  The inclusion $D\to X$ is a closed
immersion of proper $k$-schemes, so $L|_D$ is ample on $D$ by
\Cref{lem:deg-ample-basic}(3), and $\dim D=g-1$.  By the inductive hypothesis
\[
  \bigl((L|_D)^{g-1}\bigr)_D>0 .
\]

\emph{Comparison of the two Snapper polynomials.}  Put
$P(m):=\chi(X,L^{\otimes m})$ and $R(m):=\chi\bigl(D,(L|_D)^{\otimes m}\bigr)$,
both functions on $\Z$.  Tensoring the hyperplane sequence with
$L^{\otimes(m-1)}$ gives the exact sequence
\[
  0\longrightarrow L^{\otimes(m-1)}\longrightarrow L^{\otimes m}
  \longrightarrow(L|_D)^{\otimes m}\longrightarrow0,
\]
where we used $(L^{\otimes m})|_D=(L|_D)^{\otimes m}$ and the fact that the
Euler characteristic of a sheaf pushed forward along a closed immersion is
unchanged, which is the case of a closed immersion in
\Cref{prop:deg-euler-finite-morphism}(1).  By
\Cref{lem:deg-euler-additive}, $P(m)-P(m-1)=R(m)$ for every $m\in\Z$, that is
$(\Delta P)(m)=R(m+1)$.

\emph{Conclusion.}  By \Cref{lem:deg-intersection-is-top-difference} applied on
$X$ with $L_1=\cdots=L_g=L$, and by
\Cref{lem:deg-difference-explicit},
\[
  (L^g)_X=\bigl(\Delta^gP\bigr)(0)
  =\bigl(\Delta^{g-1}(\Delta P)\bigr)(0)
  =\sum_{j=0}^{g-1}(-1)^{g-1-j}\binom{g-1}{j}(\Delta P)(j)
  =\bigl(\Delta^{g-1}R\bigr)(1).
\]
By \Cref{thm:deg-snapper}, $R$ is a numerical polynomial of degree at most
$\dim D=g-1$, so $\Delta^{g-1}R$ is a constant by
\Cref{lem:deg-difference-lowers-degree}(2); by
\Cref{lem:deg-intersection-is-top-difference} on $D$ that constant is
$\bigl((L|_D)^{g-1}\bigr)_D$.  Hence
$(L^g)_X=\bigl((L|_D)^{g-1}\bigr)_D>0$.
\end{proof}

\begin{corollary}[A proper scheme with an ample trivial bundle is finite]
\label{cor:deg-ample-trivial-finite}
\uses{def:deg-ample}
Let $X$ be a proper scheme over $k$ carrying a line bundle $L$ which is at once
ample and trivial, that is $L\cong\O_X$.  Then $X$ is finite over $k$; in
particular $\dim X\leq0$ and $X$ has finitely many points.
\end{corollary}
\begin{proof}
\uses{prop:deg-ample-positive, thm:deg-snapper,
lem:deg-intersection-is-top-difference, lem:deg-difference-lowers-degree,
thm:av-zariski-main}
Let $g:=\dim X$ and suppose $g\geq1$.  The Snapper polynomial of the tuple
$(L,\dots,L)$ with $L\cong\O_X$ is the constant function
$\underline m\mapsto\chi(X,\O_X)$, so any single difference operator
annihilates it and, by \Cref{lem:deg-intersection-is-top-difference},
$(L^g)=0$.  But $L$ is ample and $\dim X=g$, so
\Cref{prop:deg-ample-positive} gives $(L^g)>0$; a contradiction.  Hence
$g\leq0$.  If $X=\emptyset$ there is nothing to prove.  Otherwise $g=0$, so
$X\to\Spec k$ has finite fibres and is of finite type, hence locally
quasi-finite, and it is proper; by
\Cref{thm:av-zariski-main} it is finite.
\end{proof}

\begin{proposition}[An abelian variety is projective]
\label{prop:deg-abelian-projective}
\uses{def:av-abelian-variety, def:deg-ample}
Every abelian variety $A$ over $k$ is projective; equivalently, by
\Cref{lem:deg-ample-basic}(4), it carries an ample line bundle.
\openobligation{AlgebraicGeometry: the projectivity of an abelian variety over
an algebraically closed field, and with it ampleness of an invertible sheaf on
a scheme.}
\end{proposition}
\begin{proof}
Not proved here.  This is \cite[\S6]{MumfordAbelianVarieties} and
\cite[I, \S7]{MilneAbelianVarieties}.  The classical argument runs as follows.
Choose an affine open $U\subseteq A$ containing the identity $e$.  Since $A$ is
smooth (\Cref{prop:av-basic}) and the complement of an affine open in a smooth
variety is of pure codimension one, the reduced complement $D:=A\setminus U$ is
an effective divisor, and $A\setminus\operatorname{Supp}D=U$ is affine.  One
then shows that $\O_A(3D)$ is very ample, by using the theorem of the square to
translate $D$ and thereby separate any pair of points and any tangent vector.
The step that makes this work, and the reason it belongs to this chapter rather
than to \Cref{chap:cube}, is that it uses only the group law and the affineness
of $U$, not the theorem of the cube.
\end{proof}

\begin{proposition}[Existence of a symmetric ample line bundle]
\label{prop:deg-symmetric-ample}
\uses{def:av-abelian-variety, def:deg-ample, def:cube-mulby-int}
Every abelian variety $A$ over $k$ carries a line bundle $L$ which is ample and
symmetric, that is $[-1]^*L\cong L$.
\end{proposition}
\begin{proof}
\uses{prop:deg-abelian-projective, lem:deg-ample-basic,
lem:cube-mulby-int-points, def:cube-mulby-int}
By \Cref{prop:deg-abelian-projective} there is an ample line bundle $L_0$ on
$A$.  By \Cref{lem:cube-mulby-int-points} the morphism $[-1]$ is the inversion
morphism of the group scheme $A$ and satisfies $[-1]\circ[-1]=[1]=\mathrm{id}_A$;
in particular $[-1]$ is an isomorphism, equal to its own inverse.  Put
\[
  L:=L_0\otimes[-1]^*L_0 .
\]
Then $[-1]^*L_0$ is ample by \Cref{lem:deg-ample-basic}(3) applied to the
isomorphism $[-1]$, and $L$ is ample by \Cref{lem:deg-ample-basic}(1) as a
tensor product of two ample bundles.  Finally, pullback is pseudofunctorial and
compatible with tensor products, so
\[
  [-1]^*L\cong[-1]^*L_0\otimes[-1]^*[-1]^*L_0
  \cong[-1]^*L_0\otimes\bigl([-1]\circ[-1]\bigr)^*L_0
  \cong[-1]^*L_0\otimes L_0\cong L,
\]
so $L$ is symmetric.
\end{proof}

\section{$[n]$ is finite, flat and surjective}\label{sec:deg-isogeny}

\begin{proposition}[Miracle flatness]\label{prop:deg-miracle-flatness}
Let $f\colon X\to Y$ be a morphism locally of finite type between locally
noetherian schemes, let $x\in X$ and $y:=f(x)$.  Suppose that $\O_{X,x}$ is
Cohen--Macaulay, that $\O_{Y,y}$ is regular and that
\[
  \dim\O_{X,x}=\dim\O_{Y,y}+\dim\O_{X_y,x},
\]
where $X_y$ is the fibre of $f$ over $y$.  Then $f$ is flat at $x$.  In
particular a finite surjective morphism between smooth varieties of the same
dimension over $k$ is flat.
\openobligation{CommutativeAlgebra: miracle flatness, that is, the local
criterion for flatness of a morphism whose source is Cohen-Macaulay and whose
target is regular and whose fibre dimension is the difference of the
dimensions. Mathlib v4.33.1 has neither the Cohen-Macaulay condition nor this
criterion.}
\end{proposition}
\begin{proof}
Not proved here.  This is the standard local criterion for flatness; see
\cite[\S23]{MatsumuraCRT} and the section \emph{Flatness and dimension} of the
chapter \emph{Commutative algebra} in \cite{stacks-project}.  For the last
sentence: a
smooth $k$-variety is regular, hence Cohen--Macaulay, and its local ring at a
closed point has dimension equal to the dimension of the variety; a finite
morphism has $0$-dimensional fibres, so the displayed equality reads
$g=g+0$ for a finite surjective morphism between $g$-dimensional smooth
varieties, and flatness at every point follows.  Flatness at the non-closed
points follows from flatness at the closed points because the closed points are
dense and flatness is an open condition for a morphism of finite type between
noetherian schemes; alternatively, apply the criterion verbatim at every point,
the two local rings there having equal dimension because both schemes are
irreducible and catenary.
\end{proof}

\begin{lemma}[The $n$-torsion is finite]\label{lem:deg-torsion-finite}
\uses{def:av-torsion, def:cube-mulby-int, def:av-abelian-variety}
Let $A$ be an abelian variety over $k$, let $n\in\Z$ with $n\neq0$ and let
$A[n]=\ker[n]$ be the $n$-torsion subgroup scheme of \Cref{def:av-torsion}.
Then $A[n]$ is finite over $k$.
\end{lemma}
\begin{proof}
\uses{prop:deg-symmetric-ample, thm:cube-mulby-pullback,
lem:deg-ample-basic, cor:deg-ample-trivial-finite, lem:av-torsion-pullback,
lem:cube-pullback-constant, prop:av-basic}
Choose a symmetric ample line bundle $L$ on $A$ by
\Cref{prop:deg-symmetric-ample}.  By \Cref{thm:cube-mulby-pullback},
\[
  [n]^*L\;\cong\;L^{\otimes n^2}.
\]

Let $j\colon A[n]\to A$ be the canonical morphism.  By
\Cref{lem:av-torsion-pullback} the square
\[
  \begin{CD}
    A[n]      @>{j}>>  A\\
    @VVV               @VV{[n]}V\\
    \Spec k   @>{e}>>  A
  \end{CD}
\]
is cartesian; in particular $j$ is a closed immersion, being the base change of
the closed immersion $e\colon\Spec k\to A$ --- a section of the separated
morphism $A\to\Spec k$ --- along $[n]$, and $A[n]$ is a closed subscheme of the
proper $k$-scheme $A$, hence proper over $k$.

Consider the line bundle $j^*[n]^*L=([n]\circ j)^*L$ on $A[n]$.  On one hand
the square commutes, so $[n]\circ j=e\circ\pi$ is the constant morphism at the
identity point $e$, where $\pi\colon A[n]\to\Spec k$ is the structure
morphism; by \Cref{lem:cube-pullback-constant} the bundle
$([n]\circ j)^*L$ is therefore \emph{trivial}.  On the other hand
$[n]^*L\cong L^{\otimes n^2}$ with $n^2\geq1$, so $[n]^*L$ is ample by
\Cref{lem:deg-ample-basic}(2), and its pullback along the closed immersion $j$
is ample by \Cref{lem:deg-ample-basic}(3).

So $A[n]$ is a proper $k$-scheme carrying a line bundle that is at once ample
and trivial.  By \Cref{cor:deg-ample-trivial-finite}, $A[n]$ is finite over
$k$.
\end{proof}

\begin{lemma}[$[n]$ is proper]\label{lem:deg-mulby-proper}
\uses{def:cube-mulby-int, def:av-abelian-variety}
Let $A$ be an abelian variety over $k$ and $n\in\Z$.  Then
$[n]\colon A\to A$ is proper.
\end{lemma}
\begin{proof}
\uses{prop:av-basic, def:cube-mulby-int}
The composite of $[n]$ with the structure morphism $A\to\Spec k$ is the
structure morphism $A\to\Spec k$, which is proper, and $A\to\Spec k$ is
separated (\Cref{prop:av-basic}).  Cancellation for proper morphisms --- if
$h\circ f$ is proper and $h$ is separated then $f$ is proper --- gives the
claim.  This is the argument already used in the proof of
\Cref{thm:av-mulby-etale-corollaries}, where it is recorded that Mathlib has it
as \texttt{AlgebraicGeometry.IsProper.of\_comp}.
\end{proof}

\begin{theorem}[$[n]$ is an isogeny for every $n\neq0$]\label{thm:deg-isogeny}
\uses{def:av-abelian-variety, def:cube-mulby-int}
Let $A$ be an abelian variety of dimension $g$ over $k$ and let $n\in\Z$,
$n\neq0$.  Then $[n]\colon A\to A$ is finite, flat and surjective.
\end{theorem}
\begin{proof}
\uses{lem:deg-torsion-finite, lem:deg-mulby-proper, thm:av-zariski-main,
lem:av-translation-spreading, prop:av-basic, prop:deg-miracle-flatness,
lem:av-torsion-pullback}
\emph{Locally quasi-finite.}  By \Cref{lem:deg-torsion-finite} the scheme
$A[n]$ is finite over $k$, and by \Cref{lem:av-torsion-pullback} it is the
fibre of $[n]$ over the identity point $e$.  So $[n]$ is quasi-finite at every
point of $A[n]$, and $A[n]\neq\emptyset$ because it contains $e$.  The locus
\[
  V:=\{x\in A:\ [n]\ \text{is quasi-finite at}\ x\}
\]
is open, quasi-finiteness at a point being an open condition for a morphism
locally of finite type between locally noetherian schemes, and it is non-empty
by the above.  Quasi-finiteness at a point is local on source and target and is
stable under pre- and post-composition with isomorphisms, so
\Cref{lem:av-translation-spreading} applies and gives $V=A$.  Thus $[n]$ is
locally quasi-finite.

\emph{Finite.}  By \Cref{lem:deg-mulby-proper} the morphism $[n]$ is proper,
and it is locally quasi-finite; by Zariski's main theorem in the form
\Cref{thm:av-zariski-main} it is finite.

\emph{Surjective.}  Being finite, $[n]$ is closed, so its set-theoretic image
$Z:=[n](A)$ is a closed subset of $A$.  It is irreducible, being the image of
the irreducible space $A$ (\Cref{prop:av-basic}).  Every fibre of $[n]$ over a
point of $Z$ is $0$-dimensional, since $[n]$ is finite, so the fibre-dimension
formula for a dominant morphism of irreducible schemes of finite type over a
field, applied to the surjection $A\to Z$ with $Z$ given its reduced structure,
gives $\dim A=\dim Z+0$, that is $\dim Z=g$.  A closed irreducible subset of
the irreducible $g$-dimensional scheme $A$ of dimension $g$ is all of $A$,
because a proper closed subset of an irreducible scheme of finite type over a
field has strictly smaller dimension.  So $Z=A$ and $[n]$ is surjective.

\emph{Flat.}  By \Cref{prop:av-basic} the abelian variety $A$ is smooth over
$k$ and irreducible of dimension $g$; so $[n]\colon A\to A$ is a finite
surjective morphism between smooth $k$-varieties of the same dimension $g$.  By
\Cref{prop:deg-miracle-flatness} it is flat.
\end{proof}

\section{The degree formula}\label{sec:deg-degree}

\begin{theorem}[The degree of $[n]$]\label{thm:deg-degree-formula}
\uses{def:av-abelian-variety, def:cube-mulby-int, def:av-degree}
Let $A$ be an abelian variety of dimension $g$ over $k$ and let $n\in\Z$,
$n\neq0$.  Then $[n]$ is finite and flat, so its degree is defined, and
\[
  \deg[n]=n^{2g}.
\]
\end{theorem}
\begin{proof}
\uses{thm:deg-isogeny, prop:deg-symmetric-ample, thm:cube-mulby-pullback,
prop:deg-projection-formula, cor:deg-intersection-powers,
prop:deg-ample-positive, prop:deg-abelian-projective, prop:av-basic,
def:av-degree}
By \Cref{thm:deg-isogeny} the morphism $[n]$ is finite, flat and surjective, so
$\deg[n]$ is defined by \Cref{def:av-degree}; $A$ is connected
(\Cref{prop:av-basic}), so the degree does not depend on the chosen point.

Choose a symmetric ample line bundle $L$ on $A$ by
\Cref{prop:deg-symmetric-ample}.  By \Cref{prop:deg-abelian-projective} the
scheme $A$ is projective, and it is integral of dimension $g$ by
\Cref{prop:av-basic}.  So \Cref{prop:deg-projection-formula} applies to the
finite surjective morphism $[n]\colon A\to A$ and to the $g$-tuple
$(L,\dots,L)$:
\[
  \bigl(([n]^*L)^g\bigr)=\deg[n]\cdot(L^g).
\]

On the other hand \Cref{thm:cube-mulby-pullback} gives
$[n]^*L\cong L^{\otimes n^2}$, because $L$ is symmetric, so by
\Cref{cor:deg-intersection-powers}
\[
  \bigl(([n]^*L)^g\bigr)=\bigl((L^{\otimes n^2})^g\bigr)
  =(n^2)^g\,(L^g)=n^{2g}\,(L^g).
\]

Comparing the two displays gives $\deg[n]\cdot(L^g)=n^{2g}\,(L^g)$.  Finally
$(L^g)>0$ by \Cref{prop:deg-ample-positive}, since $L$ is ample and
$\dim A=g$; so $(L^g)$ may be cancelled from the identity, and
$\deg[n]=n^{2g}$.
\end{proof}

\begin{remark}[What this discharges]\label{rem:deg-discharged}
\Cref{prop:deg-symmetric-ample}, \Cref{prop:deg-ample-positive},
\Cref{prop:deg-projection-formula}, \Cref{thm:deg-isogeny} and
\Cref{thm:deg-degree-formula} are the five statements that
\Cref{sec:av-cube} carried as open obligations about ampleness, intersection
numbers and the degree computation; they are
\Cref{prop:av-symmetric-ample}, \Cref{prop:av-ample-self-intersection},
\Cref{prop:av-degree-intersection}, \Cref{thm:av-isogeny-general} and
\Cref{thm:av-degree-formula}, which now cite this chapter.  What remains open
after this chapter is \emph{not} any of them but the layer beneath: the Euler
characteristic (\Cref{prop:deg-cohomology-finiteness} and
\Cref{prop:deg-euler-finite-morphism}), ampleness
(\Cref{lem:deg-ample-basic}), the projectivity of an abelian variety
(\Cref{prop:deg-abelian-projective}), the d\'evissage of a coherent sheaf
(\Cref{lem:deg-devissage}) and miracle flatness
(\Cref{prop:deg-miracle-flatness}).  Together with
\Cref{thm:cube-mulby-pullback} of \Cref{chap:cube}, this closes
\Cref{prop:av-isogeny} and the four \texttt{sorry}s listed in
\Cref{sec:av-status}, modulo that layer and modulo the descent to a general
base field of \Cref{rem:av-descent-to-any-field}.
\end{remark}

\begin{remark}[Negative $n$, and the group of $k$-points]
\label{rem:deg-negative-n}
\Cref{thm:deg-degree-formula} is stated for every non-zero integer $n$, positive
or negative, and $n^{2g}=(n^2)^g>0$ in either case; this is why
\Cref{def:cube-mulby-int}, the $\Z$-valued $[n]$, rather than
\Cref{def:av-mulby}, is the right notion here.  Nothing in the proof
distinguishes the sign: the only place $n$ enters is through $n^2$, in
$[n]^*L\cong L^{\otimes n^2}$.  In particular $[-1]$ has degree $1$, as it must,
being an isomorphism.  Note also that \Cref{thm:deg-degree-formula} counts the
degree as a rank, not as a number of points: when $n$ is divisible by the
characteristic of $k$ the scheme $A[n]$ is non-reduced and
$\#A[n](k)<n^{2g}$.  The comparison of the two counts, which is where
\Cref{prop:av-finite-etale-points} enters, is the subject of
\Cref{prop:av-etale-torsion} and not of this chapter.
\end{remark}

\section{Formalisation cost}\label{sec:deg-cost}

This section counts what Mathlib \texttt{v4.33.1} does not have and this
chapter needs.  It is written to \emph{compose} with \Cref{sec:cube-cost}, not
to repeat it: the two are meant to be added together, and nothing charged
there is charged again here.  Concretely:

\begin{itemize}
  \item The cohomology $H^i(S,F)$ of a sheaf of modules on a scheme, in all
  degrees, is notion (N4) of \Cref{sec:cube-cost} --- two definitions and about
  twelve lemmas, in
  \texttt{WeilConjectures/AlgebraicGeometry/Cohomology/Basic.lean}.  The Euler
  characteristic of this chapter sits directly on top of it and is charged
  below as (N8); $H^i$ itself is \textbf{not} recounted.
  \item Coherence of a sheaf of modules on a locally noetherian scheme is
  notion (N1) of \Cref{sec:cube-cost} --- one definition and about six lemmas.
  Every statement of this chapter is about coherent sheaves, and none of that
  is recounted.
  \item The projection formula for the direct image is charged in
  \Cref{sec:cube-cost} as part of notion (N5).
  \Cref{prop:deg-euler-finite-morphism}(2) is exactly that statement and is
  \textbf{not} recounted; only part (1) of that proposition, the invariance of
  cohomology along an affine morphism, is new and is charged below.
  \item The direct images along a proper morphism and base change for them are
  notion (N6) of \Cref{sec:cube-cost}.  This chapter uses none of that
  machinery: no semicontinuity, no Grauert, no formal functions.  Its
  cohomological needs are strictly the \emph{absolute} groups $H^i(X,F)$ over a
  field, together with the one affine-morphism statement just mentioned.
  \item The tensor product of sheaves of modules, invertible sheaves and
  $\mathrm{Pic}(X)$, and the pullback homomorphism
  $u^*\colon\mathrm{Pic}(V)\to\mathrm{Pic}(S)$, are the parallel layer of
  \Cref{warn:deg-parallel-layer} and notion (N3) of \Cref{sec:cube-cost}
  respectively; neither is recounted.
\end{itemize}

\subsection*{What is already available}

Beyond items (1) to (6) of the corresponding list in \Cref{sec:cube-cost},
which are not repeated, three further things exist and are used.
\begin{enumerate}
  \item \texttt{AlgebraicGeometry.finite\_appTop\_of\_universallyClosed} in
  \texttt{Mathlib/AlgebraicGeometry/Morphisms/Proper.lean}, which is the whole
  of \Cref{lem:deg-proper-affine-finite} once combined with the
  $\Gamma\dashv\Spec$ adjunction; and
  \texttt{AlgebraicGeometry.IsFinite.of\_isProper\_of\_locallyQuasiFinite},
  which is \Cref{thm:av-zariski-main} and is used twice here.
  \item Associated primes of a module over a noetherian ring, in
  \texttt{Mathlib/RingTheory/Ideal/AssociatedPrime/}, including their
  finiteness and the characterisation of zerodivisors.  This is the
  commutative-algebra half of \Cref{lem:deg-general-hyperplane}; the
  scheme-theoretic half is not there and is charged below.
  \item Projective space as a scheme: \texttt{AlgebraicGeometry.Proj} in
  \texttt{Mathlib/AlgebraicGeometry/ProjectiveSpectrum/} and, in this project,
  \texttt{AlgebraicGeometry.ProjectiveSpace} --- the scheme
  $\mathbf P(n;S)$ over a base scheme $S$, with its universal property under
  base change --- in
  \texttt{WeilConjectures/AlgebraicGeometry/ProjectiveSpace/Basic.lean}.  The
  twisting sheaves $\O(d)$ are also in this project now, both on $\operatorname{Proj}$ of a
  graded ring and, by pullback along the map to the integral model, on
  $\mathbf P(n;S)$, in
  \texttt{WeilConjectures/AlgebraicGeometry/TwistingSheaf.lean} and
  \texttt{WeilConjectures/AlgebraicGeometry/ProjectiveSpace/TwistingSheaf.lean}.
  What is \emph{not} there, in Mathlib or in this project, is the
  invertibility of $\O(1)$, the correspondence between morphisms to
  $\mathbf P^N$ and globally generated line bundles, or any notion of
  ampleness; those are charged below as (N11).
\end{enumerate}

\subsection*{The missing notions, one by one}

Six notions are missing beyond (N1) to (N7) of \Cref{sec:cube-cost}; they are
numbered (N8) to (N13) so that the two lists can be concatenated.  For each:
which statements of this chapter need it, how many definitions and how many
lemmas it is, and the file it would live in.  As in \Cref{sec:cube-cost}, the
lemma counts estimate the API a usable version needs, not the individual proof
steps.

\begin{enumerate}
  \item[(N8)] \emph{The Euler characteristic of a coherent sheaf on a proper
  scheme over a field.}  Needed to state \Cref{def:deg-euler}, and used by
  every statement of \Cref{sec:deg-snapper} and \Cref{sec:deg-intersection}.
  This sits directly on (N4) and (N1) of \Cref{sec:cube-cost}, which are not
  recounted.  \emph{One definition}, $\chi(X,F)=\sum_i(-1)^i\dim_kH^i(X,F)$,
  and \emph{about six lemmas}: the finite-dimensionality of $H^i(X,F)$ for $X$
  proper over a field and $F$ coherent; the vanishing $H^i=0$ for
  $i>\dim X$; the vanishing of the alternating sum of dimensions along a
  bounded exact sequence of finite-dimensional vector spaces --- the only one
  of the six that is pure linear algebra, and the only one for which
  \texttt{Mathlib/Algebra/Homology/EulerCharacteristic.lean} is a plausible
  starting point; additivity of $\chi$ on short exact sequences,
  \Cref{lem:deg-euler-additive}; $\chi(X,F)=\dim_k\Gamma(X,F)$ for $X$ affine,
  \Cref{lem:deg-euler-affine}; and invariance of $H^i$, hence of $\chi$, along
  an affine morphism, which is \Cref{prop:deg-euler-finite-morphism}(1).  In a
  new file
  \texttt{WeilConjectures/AlgebraicGeometry/Cohomology/EulerCharacteristic.lean},
  beside the \texttt{Basic.lean} of (N4).  The first two of the six are the
  expensive ones and are the reason this notion cannot be cut loose from
  (N4): they are the finiteness theorem for the cohomology of a proper
  morphism, in the absolute case.
  \item[(N9)] \emph{Numerical polynomials in several variables and the
  difference calculus.}  Needed to state \Cref{def:deg-numerical} and
  \Cref{thm:deg-snapper}, and used by every statement of
  \Cref{sec:deg-intersection} and by \Cref{prop:deg-ample-positive}.
  \emph{Two definitions} --- the predicate ``$P\colon\Z^r\to\Z$ is a numerical
  polynomial of total degree at most $d$'' and the difference operator
  $\Delta_v$ --- and \emph{about seven lemmas}, all four of the section proved
  in full here (\Cref{lem:deg-numerical-unique},
  \Cref{lem:deg-difference-lowers-degree},
  \Cref{lem:deg-difference-criterion} and
  \Cref{lem:deg-difference-explicit}) plus three pieces of API they need: the
  binomial coefficient $\binom{m}{i}$ as an integer-valued polynomial on $\Z$,
  the Pascal identity for it, and the substitution of linear forms into a
  polynomial of bounded total degree, which is what makes the reduction to
  very ample bundles work inside \Cref{thm:deg-snapper}.  In a new file
  \texttt{WeilConjectures/RingTheory/NumericalPolynomial.lean}.
  \emph{This notion is the one item of the whole chapter that can be
  formalised today}: it mentions no scheme, no sheaf and no cohomology, it
  needs nothing from \Cref{sec:cube-cost}, and it is about a fifth of the
  chapter by length.  Mathlib \texttt{v4.33.1} has
  \texttt{Polynomial.hilbertPoly} in
  \texttt{Mathlib/RingTheory/Polynomial/HilbertPoly.lean}, but that is the
  Hilbert polynomial of a graded module in \emph{one} variable and over a
  field of characteristic zero, and it is not the multivariable
  integer-valued statement needed here.
  \item[(N10)] \emph{Intersection numbers of line bundles on a projective
  scheme.}  Needed to state \Cref{def:deg-intersection},
  \Cref{prop:deg-projection-formula} and \Cref{prop:deg-ample-positive}, and
  hence \Cref{thm:deg-degree-formula}.  \emph{Two definitions} --- the
  intersection number $(L_1\cdots L_d\cdot F)$ with coefficients in a coherent
  sheaf and the abbreviation $(L^d)$ for the top self-intersection --- and
  \emph{about eight lemmas}: the identification with the top difference of the
  Snapper polynomial (\Cref{lem:deg-intersection-is-top-difference}); symmetry
  (\Cref{lem:deg-intersection-symmetric}); additivity in the coefficient sheaf
  (\Cref{lem:deg-intersection-additive}); vanishing below the dimension
  (\Cref{lem:deg-intersection-small-support}); the three parts of
  multilinearity (\Cref{thm:deg-intersection-multilinear}); and the scaling
  rule (\Cref{cor:deg-intersection-powers}).  In a new file
  \texttt{WeilConjectures/AlgebraicGeometry/IntersectionNumber.lean}.
  Mathlib \texttt{v4.33.1} has
  \texttt{Mathlib/AlgebraicGeometry/AlgebraicCycle/Basic.lean} and
  \texttt{Mathlib/AlgebraicGeometry/OrderOfVanishing.lean}, which are the
  beginnings of a \emph{cycle-theoretic} intersection theory; that is a
  different route from Snapper's and none of it applies directly.
  \item[(N11)] \emph{Ampleness and very ampleness of an invertible sheaf, and
  global generation.}  Needed to state \Cref{def:deg-ample},
  \Cref{lem:deg-ample-basic}, \Cref{prop:deg-ample-positive},
  \Cref{cor:deg-ample-trivial-finite}, \Cref{prop:deg-abelian-projective} and
  \Cref{prop:deg-symmetric-ample}, and used by \Cref{lem:deg-torsion-finite},
  \Cref{thm:deg-isogeny} and \Cref{thm:deg-degree-formula}.  \emph{Three
  definitions} --- a sheaf of modules generated by its global sections,
  ampleness, and very ampleness relative to a base --- and \emph{about nine
  lemmas}: the twisting sheaf $\O(1)$ on $\mathbf P^N$ and its ampleness; the
  correspondence between a morphism to $\mathbf P^N$ and a globally generated
  line bundle with $N+1$ chosen sections, which is what makes
  \Cref{lem:deg-general-hyperplane} available; the fact that an immersion of a
  proper scheme is a closed immersion, which is what makes
  \Cref{def:deg-ample} usable; and the five permanence properties collected as
  \Cref{lem:deg-ample-basic}, of which the fourth --- ample if and only if some
  power is very ample --- is much the hardest and is a genuine theorem, not
  API.  In two new files
  \texttt{WeilConjectures/AlgebraicGeometry/Ample.lean} and
  \texttt{WeilConjectures/AlgebraicGeometry/TwistingSheaf.lean}.
  \emph{This is the largest single item in this chapter's count}, and the one
  most obviously worth upstreaming: nothing in it mentions abelian varieties.
  \item[(N12)] \emph{Associated points of a coherent sheaf on a scheme, the
  d\'evissage of a coherent sheaf, and the generic rank.}  Needed to state
  \Cref{lem:deg-general-hyperplane} and \Cref{lem:deg-devissage}, and used by
  \Cref{thm:deg-snapper}, \Cref{prop:deg-intersection-rank},
  \Cref{prop:deg-projection-formula} and \Cref{prop:deg-ample-positive}.
  \emph{Two definitions} --- the set of associated points of a coherent sheaf
  and the generic rank of a coherent sheaf on an integral scheme --- and
  \emph{about six lemmas}: finiteness of the set of associated points; the
  fact that every generic point of an irreducible component of the support is
  associated; the criterion for a section of a line bundle to be a
  nonzerodivisor on a coherent sheaf; the d\'evissage
  \Cref{lem:deg-devissage}; that a proper closed subset of an irreducible
  scheme of finite type over a field has strictly smaller dimension; and
  Krull's principal ideal theorem in the scheme-theoretic form used in
  \Cref{prop:deg-ample-positive}.  In
  \texttt{WeilConjectures/Mathlib/AlgebraicGeometry/Noetherian.lean} for the
  first four and
  \texttt{WeilConjectures/Mathlib/AlgebraicGeometry/Properties.lean} for the
  last two, both mirroring existing Mathlib files.  The ring-theoretic half is
  available (see item (2) above); what is missing is entirely the passage to
  schemes.
  \item[(N13)] \emph{Miracle flatness.}  Needed to state
  \Cref{prop:deg-miracle-flatness} and used only by the flatness clause of
  \Cref{thm:deg-isogeny}.  \emph{One definition}, the Cohen--Macaulay
  condition on a noetherian local ring --- which Mathlib \texttt{v4.33.1} does
  not have in any form; \texttt{grep} for \texttt{CohenMacaulay} over the whole
  library returns nothing --- and \emph{about four lemmas}: a regular local
  ring is Cohen--Macaulay; the local criterion for flatness in the form of
  \Cref{prop:deg-miracle-flatness}; the fibre-dimension formula
  $\dim\O_{X,x}=\dim\O_{Y,y}+\dim\O_{X_y,x}$ for a finite morphism of
  irreducible varieties of the same dimension; and the transfer of the local
  criterion to the scheme level.  In
  \texttt{WeilConjectures/Mathlib/RingTheory/Flat/Basic.lean} and
  \texttt{WeilConjectures/Mathlib/AlgebraicGeometry/Morphisms/Flat.lean}.
  This is the one notion of the chapter that is pure commutative algebra and
  needs no cohomology; it is also the only one whose absence was already
  recorded in \Cref{chap:av-isogeny}, in the open obligation on
  \Cref{thm:av-isogeny-general}.
\end{enumerate}

Adding the columns: \emph{eleven definitions} and \emph{about forty lemmas of
supporting API}, on top of which the statements of this chapter sit.  One
further definition belongs to no notion above --- the predicate
``$L$ is symmetric'' on a commutative group scheme, part of
\Cref{def:deg-ample}, which is glue over the $\Z$-valued $[n]$ that the
project already has --- and it brings the definition count to twelve.

\subsection*{Definitions missing: twelve, of which ten carry content}

\begin{enumerate}
  \item \emph{The Euler characteristic $\chi(X,F)$}, from (N8).
  \item \emph{Numerical polynomial of total degree at most $d$}, from (N9).
  \item \emph{The difference operator $\Delta_v$}, from (N9).
  \item \emph{The intersection number $(L_1\cdots L_d\cdot F)$}, from (N10).
  \item \emph{The top self-intersection $(L^d)$}, from (N10).  Glue: the
  special case of the previous one with all entries equal.
  \item \emph{Generated by global sections}, from (N11).
  \item \emph{Ample}, from (N11).
  \item \emph{Very ample}, from (N11).
  \item \emph{The associated points of a coherent sheaf}, from (N12).
  \item \emph{The generic rank of a coherent sheaf}, from (N12).
  \item \emph{Cohen--Macaulay local ring}, from (N13).
  \item \emph{Symmetric line bundle on a commutative group scheme}, part of
  \Cref{def:deg-ample}.  Glue: $[-1]^*L\cong L$ for the $[-1]$ of
  \Cref{def:cube-mulby-int}.
\end{enumerate}
Items (1) to (4) and (6) to (11) carry content; items (5) and (12) are
abbreviations.  None of the twelve duplicates any of the ten counted in
\Cref{sec:cube-cost}, so the two chapters together need
\emph{twenty-two definitions} and \emph{about eighty-five lemmas of supporting
API}.

\subsection*{Statements missing: six left open, twenty-two proved here}

Six statements are left open, and they are the only unproved statements in the
chapter.
\begin{enumerate}
  \item \Cref{prop:deg-cohomology-finiteness} needs (N8), and beneath it (N4)
  and (N1) of \Cref{sec:cube-cost}: finiteness and vanishing of the cohomology
  of a coherent sheaf on a projective scheme over a field.  Used by every
  statement of the chapter.
  \item \Cref{prop:deg-euler-finite-morphism} needs (N8) for part (1) and
  (N5) of \Cref{sec:cube-cost} for part (2).  Used by
  \Cref{prop:deg-projection-formula} and \Cref{prop:deg-ample-positive}.
  \item \Cref{lem:deg-ample-basic} needs (N11).  Used by
  \Cref{thm:deg-snapper}, \Cref{prop:deg-ample-positive},
  \Cref{lem:deg-torsion-finite} and \Cref{prop:deg-symmetric-ample}.
  \item \Cref{lem:deg-devissage} needs (N12).  Used by
  \Cref{prop:deg-intersection-rank} only.
  \item \Cref{prop:deg-abelian-projective} needs (N11) and the classical
  construction of a very ample divisor on an abelian variety.  Used by
  \Cref{prop:deg-symmetric-ample} and \Cref{thm:deg-degree-formula}.  This is
  the one open statement of the chapter that is about abelian varieties rather
  than about general algebraic geometry, and the only one that will not be
  supplied by a general-purpose Lean subtree.
  \item \Cref{prop:deg-miracle-flatness} needs (N13).  Used by the flatness
  clause of \Cref{thm:deg-isogeny} only.
\end{enumerate}

Twenty-two statements are proved here in full and are new to Lean:
\Cref{def:deg-euler} and \Cref{lem:deg-euler-additive},
\Cref{lem:deg-euler-affine},
\Cref{lem:deg-numerical-unique},
\Cref{lem:deg-difference-lowers-degree},
\Cref{lem:deg-difference-criterion},
\Cref{lem:deg-difference-explicit},
\Cref{lem:deg-proper-affine-finite},
\Cref{lem:deg-general-hyperplane},
\Cref{thm:deg-snapper},
\Cref{lem:deg-intersection-is-top-difference},
\Cref{lem:deg-intersection-symmetric},
\Cref{lem:deg-intersection-additive},
\Cref{lem:deg-intersection-small-support},
\Cref{thm:deg-intersection-multilinear},
\Cref{cor:deg-intersection-powers},
\Cref{prop:deg-intersection-rank},
\Cref{prop:deg-projection-formula},
\Cref{prop:deg-ample-positive},
\Cref{cor:deg-ample-trivial-finite},
\Cref{prop:deg-symmetric-ample},
\Cref{lem:deg-torsion-finite},
\Cref{lem:deg-mulby-proper},
\Cref{thm:deg-isogeny} and
\Cref{thm:deg-degree-formula} --- twenty-five items, of which
\Cref{def:deg-euler} is a definition and \Cref{lem:deg-mulby-proper} and
\Cref{lem:deg-euler-affine} are one-line consequences of facts Mathlib already
has, leaving twenty-two of real formalisation work.  Four are theorems, two are
corollaries, four are propositions and the remaining twelve are supporting
lemmas.

Four of the twenty-two use nothing that is missing beyond the definitions
listed above and can therefore be formalised today:
\Cref{lem:deg-numerical-unique},
\Cref{lem:deg-difference-lowers-degree},
\Cref{lem:deg-difference-criterion} and
\Cref{lem:deg-difference-explicit} --- the whole of
\Cref{sec:deg-numerical}, which is notion (N9) and is self-contained
combinatorics.  A fifth, \Cref{lem:deg-proper-affine-finite}, needs only what
Mathlib already has.  None of the others is formalisable before (N8) and (N11)
exist.

\subsection*{Summary, and how it composes with \Cref{sec:cube-cost}}

Twelve definitions, of which ten carry content; about forty lemmas of
supporting API in the six missing notions (N8) to (N13); six statements left
open; and twenty-two statements proved in full here, five of which are
formalisable immediately.

Taken together with \Cref{sec:cube-cost} --- ten definitions, about forty-five
lemmas, seven statements left open, twenty-five proved --- the two chapters
that stand between \Cref{prop:av-isogeny} and a complete Lean proof of
$\deg[n]=n^{2g}$ come to \emph{twenty-two definitions}, \emph{about
eighty-five lemmas of supporting API in thirteen missing notions}, \emph{thirteen
statements left open} and \emph{forty-seven statements proved in the
blueprint}.

The thirteen open statements are not thirteen independent theories.  They rest
on four:
\begin{enumerate}
  \item \emph{The cohomology of coherent sheaves on a scheme, with proper base
  change} --- (N1), (N4), (N5), (N6) of \Cref{sec:cube-cost} and (N8) here.
  Eight of the thirteen open statements rest on this --- six of the seven of
  \Cref{sec:cube-cost} and both of the cohomological ones here --- and it is by
  a wide margin the largest single cost of the two chapters.
  \Cref{chap:cube} needs degrees zero and one in a family; this chapter needs
  all degrees but only absolutely,
  over a field, with no base.  The two demands are close enough that one
  development serves both, but neither is a special case of the other.
  \item \emph{Ampleness} --- (N11) here.  Two open statements,
  \Cref{lem:deg-ample-basic} and \Cref{prop:deg-abelian-projective}, and no
  overlap with \Cref{chap:cube} at all: the theorem of the cube never mentions
  an ample bundle.
  \item \emph{Coherent-sheaf bookkeeping on a noetherian scheme} --- (N12)
  here, together with (N7) of \Cref{sec:cube-cost}.  Two open statements,
  \Cref{lem:deg-devissage} and
  \Cref{prop:cube-global-sections-flat-base-change}; cheap, and needing no
  cohomology.
  \item \emph{Miracle flatness} --- (N13) here.  One open statement, one
  clause of one theorem, and pure commutative algebra.
\end{enumerate}

Two practical consequences for the Lean split.  First, notion (N9), the
numerical-polynomial calculus, is a genuinely independent piece of about seven
lemmas which can be formalised now, before any decision about the cohomology
layer, and which nothing else in the project blocks; it should be cut out
first.  Second, notion (N11), ampleness, is the largest item here and is
entirely general algebraic geometry: it depends on none of the cohomological
notions, it would be a natural Mathlib contribution, and it is the gate through
which \Cref{prop:deg-ample-positive}, \Cref{lem:deg-torsion-finite} and hence
both main theorems of this chapter pass.  The three-way split ``numerical
polynomials'', ``ampleness'' and ``Euler characteristic and intersection
numbers'' is the natural one, in that order of increasing cost, and only the
third of the three is blocked on \Cref{sec:cube-cost}.
